% =====================================================================
%  capitulos/cap03.tex — Capítulo 3. El sonido y la estructura fonológica
%  Autora: Mercedes Güemes
%  Fuente: Capitulo_03_normalizado.docx
%  Transcripciones: \fon{...} = [ ], \fonol{...} = / /
% =====================================================================
\capitulo{El sonido y la estructura fonológica}{Mercedes Güemes}

\section{Introducción}

En este capítulo se abordan las principales características de los sonidos y
su valor en los sistemas lingüísticos. En las primeras secciones se revisa el
sistema ortográfico del español, considerando el principio fonémico en que se
basa y sus irregularidades. Luego, se presenta el sistema de notación del
Alfabeto Fonético Internacional (AFI), como herramienta para la descripción
articulatoria de los sonidos vocálicos y consonánticos. Posteriormente, el
capítulo aborda nociones centrales de la teoría fonológica, como la distinción
entre fono y fonema, la distribución contrastiva, complementaria y la variación
libre. Asimismo, se describen dos de los principales procesos fonológicos del
español: la asimilación y la neutralización. También se introduce el concepto
de sílaba y se presenta la configuración de su estructura interna. Por último,
se explican los rasgos suprasegmentales más relevantes del español: el acento y
la entonación.

\section{Dos disciplinas complementarias: fonética y fonología}

Las lenguas están formadas por un número limitado de unidades fónicas sin
significado, cuyo contraste permite distinguir un número ilimitado de
significados. Estas unidades, conocidas como \termino{fonemas}, deben
combinarse entre sí para poder construir unidades significativas mayores, como
por ejemplo los morfemas,\footnote{Este concepto se presenta y desarrolla en el
\capref{4}.} las palabras, los enunciados y los textos. Como manifiestan
\textcite[12]{casado2011}, «el conjunto potencialmente infinito de expresiones
que se pueden producir en una lengua está formado por un inventario limitado de
unidades sonoras discretas». A pesar de que en la actualidad buena parte de las
interacciones lingüísticas se basan en la escritura, el lenguaje humano es
fundamentalmente un fenómeno de naturaleza oral.\footnote{La mayor parte de las
lenguas humanas son orales. No obstante, también existen las lenguas de señas;
su gramática y léxico se manifiestan mediante la producción de gestos que se
perciben visualmente.} Las dos disciplinas que se encargan del estudio del
sonido humano y su función comunicativa son la \termino{fonética} y la
\termino{fonología}.

La \termino{fonética} se ocupa de estudiar el aspecto material de los sonidos
del habla, es decir, los que efectivamente pronunciamos, y se interesa
principalmente en tres aspectos que definen las tres ramas que la conforman,
según se muestra en la tabla~\ref{tab:03-ramas}.

\begin{table}[htbp]
  \centering
  \begin{tabularx}{\linewidth}{@{}L L L@{}}
    \toprule
    \textbf{Fonética acústica} & \textbf{Fonética articulatoria} & \textbf{Fonética auditiva} \\
    \midrule
    Se enfoca en la estructura de los sonidos y grupos de sonidos. Se
    concentra principalmente en las características físicas de la onda sonora.
    & Analiza cómo se producen y se articulan los sonidos en el aparato
    fonatorio humano e investiga la fisiología y la anatomía humana vinculada
    con la producción del sonido. También clasifica los sonidos de acuerdo con
    los parámetros que los generan.
    & Estudia la manera en la que el oído recibe y percibe los sonidos. Su
    objetivo es entender el funcionamiento psicoacústico del sistema auditivo y
    determinar el alcance de la audición en general. \\
    \bottomrule
  \end{tabularx}
  \caption{Ramas de la fonética}\label{tab:03-ramas}
\end{table}

La \termino{fonología}, por otra parte, es la subdisciplina de la lingüística
encargada de identificar y describir el comportamiento sistemático de los
sonidos de las lenguas. En consecuencia, para la fonología no es del todo
relevante la materialidad del sonido, sino, según iremos viendo a lo largo de
este capítulo, la manera en que los rasgos sonoros se configuran en un sistema
fonológico particular para contrastar significados. La fonología se encarga de
estudiar el conjunto abstracto de representaciones sonoras compartido por una
comunidad de hablantes y las variaciones recurrentes de los sonidos de habla
que se dan dentro de cada variedad lingüística particular.

\section{Segmentación y representación grafémica}\label{sec:03-grafemas}

Dado que las lenguas son entidades principalmente orales, el ser humano cuenta
con la escritura como forma de representar y almacenar las palabras que se
emiten. Para poder hacer un estudio sobre los sonidos de las lenguas, es
importante comprender las formas en las que la escritura alfabética transcribe
los sonidos lingüísticos. Actualmente, una gran mayoría de las lenguas del
mundo cuenta con un sistema de escritura que se basa en el principio fonémico;
esto significa que la ortografía intenta representar los sonidos del habla
asignando para cada sonido una letra. Sin embargo, los abecedarios no cumplen
con exactitud este principio; por ejemplo, algunas letras (\termino{grafemas})
no representan un único sonido: hay sonidos que se transcriben mediante dos o
más grafemas, y hay grafemas que no se pronuncian. Dicho de otra manera, la
correspondencia fonema-grafema no es perfecta. Por ejemplo, la palabra
\forma{higuera} presenta un grafema que no tiene manifestación sonora, como la
\forma{h}, y la dupla \forma{gu} se pronuncia como un solo sonido, equivalente
al primero de la palabra \forma{gala}.

Como puede observarse, la escritura alfabética presenta particularidades que la
alejan, en mayor o menor medida, de la realidad oral. La distancia que se
produce entre la oralidad y la escritura varía según las lenguas. En comparación
con otras lenguas, el español cuenta con un sistema ortográfico mayormente
transparente, dado que las irregularidades ortográficas no son tantas y que
muchas pueden resumirse en un conjunto de reglas predecibles. Sin embargo,
lenguas como el inglés o el francés\footnote{La ortografía del inglés y el
francés no es transparente debido a su evolución histórica, las influencias de
otros idiomas y la falta de reformas ortográficas que mantuvieran la ortografía
al día con los cambios que se fueron produciendo en la pronunciación.} cuentan
con sistemas ortográficos altamente opacos: no es suficiente conocer cómo se
pronuncia una palabra para saber cómo se escribe y tampoco es muy claro cómo
debe leerse una secuencia de letras. Por este motivo, para poder estudiar los
fenómenos vinculados a la oralidad, es muy importante poder tomar distancia de
la representación alfabética del sonido. Esta no es tarea fácil ya que, en el
proceso de alfabetización se producen cambios estructurales en el cerebro que
generan que nuestras actividades perceptuales de escucha estén vinculadas a la
escritura \parencite[253]{dehaene2014}. Para lograr entrar en los estudios de
los sonidos del habla, este capítulo se propone como paso inicial identificar
los límites de la escritura alfabética del español para, posteriormente,
introducir un sistema de representación fonética más adecuado que permita
reconocer, estudiar y contrastar los sonidos lingüísticos a partir de
principios más generales.

\begin{ejercicio}[\textit{El viento}, de Manu Chao]\label{ej:03-viento}
A continuación se presenta la canción de Manu Chao \textit{El viento}. Luego, a
partir de la letra y del audio, proponemos realizar una serie de ejercicios. El
audio puede escucharse en \url{https://www.youtube.com/watch?v=OK-j4WpClMs}. En
cada caso, se muestra como guía la resolución del primer ítem.

% La letra de la canción se incorpora a mano desde el original (texto protegido).
\textoaincorporar{la letra de \textit{El viento} (Manu Chao)}

\begin{consignas}
  \item Indique cuántos segmentos del habla oral se producen en cada una de las
        siguientes palabras:
        \begin{exe}
          \ex \forma{Ruta}: \respuesta{4}
          \ex \forma{Por}:
          \ex \forma{Va}:
          \ex \forma{Manu Chao}:
          \ex \forma{Carretera}:
          \ex \forma{Hombre}:
        \end{exe}
        \begin{resuelto}
        1. 4. Todos los grafemas de la palabra representan un sonido.
        \end{resuelto}
  \item Escuche la canción y observe la pronunciación de los grafemas marcados
        en negrita. ¿Cómo se puede describir la relación de las letras con los
        sonidos que representan?
        \begin{exe}
          \ex \forma{\textbf{B}abylon} -- \forma{\textbf{v}a}
          \ex \forma{\textbf{R}uta} -- \forma{ca\textbf{rr}etera} -- \forma{volve\textbf{r}á}
          \ex \forma{\textbf{S}in} -- \forma{ra\textbf{z}ón}
          \ex \forma{\textbf{H}ombre} -- \forma{\textbf{h}ambre}
        \end{exe}
        \begin{resuelto}
        1. La pronunciación de los grafemas \forma{b} y \forma{v} no presenta
        diferencias en español (con excepción de algunos hablantes que los
        distinguen a causa del contacto con otras lenguas en las que sí hay
        diferencia, como ocurre en francés o en catalán).
        \end{resuelto}
  \item El grafema inicial de las palabras \forma{cielo} y \forma{celos}
        comparte características de pronunciación con los siguientes pares,
        ¿qué ocurre en cada caso? Marque la respuesta correcta.

        \begin{tabular}{@{}l Y{2.6cm} Y{2.6cm}@{}}
          \toprule
          & Comparten grafema & Comparten fonema \\
          \midrule
          \forma{\textbf{s}in} -- \forma{ra\textbf{z}ón} & & \\
          \forma{\textbf{c}arretera} -- \forma{\textbf{c}uando} & & \\
          \bottomrule
        \end{tabular}
  \item A partir de lo visto en los puntos a) y b), ¿qué conclusiones son
        verdaderas sobre la relación grafema-fonema en español? Justifique cada
        respuesta y aporte ejemplos ilustrativos de cada caso.

        \begin{tabularx}{\linewidth}{@{}l L Y{.5cm} Y{.5cm} L@{}}
          \toprule
          & & V & F & Ejemplo \\
          \midrule
          1. & Dos o más grafemas pueden representar un único sonido. & X & &
               Los grafemas \forma{s}, \forma{c} y \forma{z} se pronuncian igual
               en comunidades seseantes. \\
          2. & Un único grafema puede representar varios sonidos. & & & \\
          3. & Cada grafema del español representa un único sonido. & & & \\
          4. & Dos grafemas contiguos pueden corresponderse con un único sonido. & & & \\
          5. & Todos los grafemas representan un sonido. & & & \\
          6. & Según el contexto, algunos grafemas cambian su pronunciación. & & & \\
          \bottomrule
        \end{tabularx}
  \item De las afirmaciones verdaderas encontradas en el punto anterior,
        identifique cuál ocurre con los grafemas señalados en estos grupos de
        palabras:
        \begin{exe}
          \ex \forma{\textbf{g}ente -- \textbf{j}efe -- \textbf{j}irafa -- \textbf{g}irar}
          \ex \forma{\textbf{g}uante -- \textbf{g}ala -- \textbf{g}erminar -- \textbf{g}ol -- \textbf{g}itano}
          \ex \forma{ta\textbf{x}i -- e\textbf{x}amen -- lé\textbf{x}ico -- ne\textbf{x}o}
          \ex \forma{\textbf{ll}amar -- \textbf{y}a -- \textbf{y}erba -- \textbf{ll}eno}
          \ex \forma{\textbf{g}anar -- \textbf{gu}eto -- \textbf{g}ol -- \textbf{gu}iño -- \textbf{g}ula}
          \ex \forma{\textbf{c}osa -- \textbf{c}álido -- \textbf{qu}echua -- \textbf{qu}ímica -- \textbf{c}ubo}
        \end{exe}
        \begin{resuelto}
        1. Dos o más grafemas pueden representar un único sonido. En este caso,
        \forma{g} (delante de vocales anteriores \fon{e} \fon{i}) y \forma{j}
        representan el mismo fonema \fonol{x}.
        \end{resuelto}
\end{consignas}
\end{ejercicio}

\begin{recuadro}[Seseo y yeísmo]\label{rec:03-seseo}
En algunas regiones hispanohablantes las palabras \forma{caza} y \forma{casa}
se distinguen fonológicamente: \fonol{ˈkaθa} y \fonol{ˈkasa} respectivamente.
Sin embargo, para los hablantes de otras regiones, a las que se denomina
seseantes, existe una única pronunciación para el fonema \fonol{s}, ya que el
fonema \fonol{θ} no se registra, tal como ocurre por ejemplo en el español de
Argentina. Se denomina \termino{seseo}, entonces, a la no diferenciación entre
los sonidos \fonol{s} y \fonol{θ} de \forma{caza} y \forma{casa}. A nivel
ortográfico, las comunidades seseantes representan el fonema \fonol{s} con los
grafemas \forma{c}, \forma{z} o \forma{s}, como ocurre en el primer segmento de
\forma{sapo} \fonol{ˈsapo}, \forma{zapato} \fonol{saˈpato} y \forma{cima}
\fonol{ˈsima}. Por otra parte, muchas comunidades hispanohablantes no presentan
el contraste entre las palabras \forma{calló} y \forma{cayó}. Esto se conoce
como \termino{yeísmo}. La fusión de ambos fonemas, en general, se produce con el
sonido palatal \fon{ʝ}. Sin embargo, en el Río de la Plata, se articula con los
sonidos postalveolares \fon{ʒ} y \fon{ʃ}, fenómeno que se conoce como
\termino{yeísmo rehilado} (v. \capref{15}).
\end{recuadro}

\begin{ejercicio}\label{ej:03-ortografia}
\begin{consignas}
  \item Las siguientes imágenes presentan errores que se vinculan con las
        irregularidades del sistema ortográfico del español. Encuentre los
        errores y explique, a partir de la relación grafema-fonema, por qué se
        han producido. Se muestra la resolución del primer ejemplo como guía.

        \imagenpendiente{Imágenes del ejercicio (carteles con errores
        ortográficos)}

        \begin{exe}
          \ex La ortografía correcta para la palabra \forma{targetas} es
              \respuesta{tarjetas}.
          \ex La ortografía correcta para la palabra \forma{benta} es\ldots
          \ex La ortografía correcta para la palabra \forma{ielo} es\ldots
          \ex La ortografía correcta para la palabra \forma{fabor} es\ldots
          \ex La ortografía correcta para la palabra \forma{cilencio} es\ldots
          \ex La ortografía correcta para la palabra \forma{descancen} es\ldots
          \ex La ortografía correcta para la palabra \forma{tullo} es\ldots
          \ex La ortografía correcta para la palabra \forma{poyo} es\ldots
        \end{exe}
        \begin{resuelto}
        1. La ortografía correcta es \forma{tarjetas}. La confusión se produce
        porque los grafemas \forma{j} y \forma{g} delante de las vocales
        \fon{e}/\fon{i} tienen una pronunciación idéntica.
        \end{resuelto}
  \item Busque ejemplos que evidencien confusiones ortográficas recurrentes
        como el caso que se produce entre \forma{haya} (árbol; verbo
        \forma{haber}) y \forma{halla} (verbo \forma{hallar}).
\end{consignas}
\end{ejercicio}

\section{Producción del sonido}

Para comprender las características básicas de los sonidos, primero es
fundamental visualizar cómo se producen y cuáles son los factores que permiten
lograr la configuración particular de cada uno de ellos. A nivel evolutivo, la
especie humana no ha desarrollado un sistema o aparato específico para producir
sonidos lingüísticos, sino que ha utilizado algunos órganos del aparato
respiratorio (pulmones, laringe, faringe, nariz) y del aparato masticador (boca,
dientes, alvéolos) para la función del habla. Así conformado, el aparato
fonatorio humano es un sistema que permite generar energía sonora a partir de
diferentes mecanismos. La respiración pulmonar es uno de esos mecanismos: el
aire procede de los pulmones y se dirige hacia el exterior atravesando la
cavidad oral o nasal para, de esta forma, emitir sonidos.

La creación de sonido es posible gracias a una serie de presiones y movimientos
que va experimentando el aire en su salida. En el trayecto del pulmón a las
cavidades de salida, boca y nariz (ver figura~\ref{fig:03-aparato}), el aire
pasa por la laringe hacia la faringe. En la laringe, la corriente de aire se
encuentra con los pliegues o cuerdas vocales que forman una abertura llamada
\termino{glotis}. El grado de abertura de la glotis varía según el sonido que se
produzca. Cuando respiramos sin emitir sonido o durante la producción de algunas
consonantes, la glotis está abierta (ver figura~\ref{fig:03-glotis}A). Sin
embargo, para que se configuren los distintos sonidos, las cuerdas deben tensarse
o vibrar de diversas maneras; esto genera que el espacio glotal varíe o se cierre
cuando hablamos (ver figura~\ref{fig:03-glotis}B). A este mecanismo se lo
denomina \termino{fonación}.

\figurapendiente{Aparato fonatorio\protect\pendiente{Falta la figura en el original}}{fig:03-aparato}

\figurapendiente{Salida del aire por la boca y por la nariz\protect\pendiente{Falta la figura en el original}}{fig:03-velo}

\figurapendiente{Glotis y cuerdas vocales. A la izquierda (A), glotis abierta;
a la derecha (B), glotis cerrada}{fig:03-glotis}

Una vez que el aire pasa por la glotis, se encuentra con una serie de cavidades,
comúnmente denominadas supraglóticas: la faringe, la cavidad oral y la cavidad
nasal. Su función principal es generar resonancia para que los sonidos se
modulen; por ello se los denomina órganos resonadores. Los sonidos
faríngeos\footnote{Los sonidos consonánticos faríngeos se producen cuando el
punto de articulación involucra la faringe. Pueden ser oclusivos o fricativos.
Si bien su presencia es escasa en español, son principalmente abundantes en
lenguas semíticas.} se producen por el golpe de aire contra la pared posterior
de la faringe. Los sonidos nasales, como \fon{m} y \fon{n}, se producen cuando el
velo del paladar permanece bajo y permite el egreso del aire a través de la nariz
(ver figura~\ref{fig:03-velo}). Los sonidos orales, por ejemplo \fon{p} y
\fon{s}, se producen cuando el aire sale por la boca; esto es posible cuando el
velo del paladar asciende y bloquea el paso del aire hacia la nariz.

La mayor cantidad de sonidos del habla se produce gracias a la complejidad que
presenta la cavidad bucal humana, integrada por órganos móviles como los labios,
la lengua (con sus partes ápice, dorso y raíz), el paladar blando, el velo del
paladar y la úvula, denominada popularmente \forma{campanilla}. También la boca
contiene puntos fijos de articulación: los dientes, las encías (alvéolos) y el
paladar duro. La diversidad de sonidos que encontramos en las lenguas va a estar
determinada por la acción de las cuerdas vocales en la glotis (si vibran o no),
el punto de contacto de los órganos móviles y fijos de la boca, y cómo se regula
la salida de aire durante todo el proceso. En la próxima sección, se explicarán
los rasgos que permiten clasificar los sonidos de las lenguas humanas.

\begin{ejercicio}\label{ej:03-aparato}
\begin{consignas}
  \item ¿Cuál es el orden correcto en el que sale el aire desde los pulmones
        hacia la boca?
        \begin{exe}
          \ex Pulmones, tráquea, faringe, laringe, cavidad bucal.
          \ex Pulmones, tráquea, cavidad bucal, laringe, faringe.
          \ex Pulmones, tráquea, laringe, faringe, cavidad bucal.
        \end{exe}
        \begin{resuelto}[Punto (a) resuelto]
        El orden correcto es 3.
        \end{resuelto}
  \item Si el velo del paladar está bajo, el sonido que se produce es:
        \begin{exe}
          \ex Oral.
          \ex Nasal.
          \ex Faríngeo.
        \end{exe}
  \item ¿Qué es la \termino{glotis}?
        \begin{exe}
          \ex Un órgano encargado de hacer vibrar las cuerdas vocales.
          \ex Un músculo que produce la fonación.
          \ex El espacio que se forma entre las cuerdas vocales.
        \end{exe}
  \item Distinga las siguientes partes de la boca en órganos fijos y móviles:

        lengua -- alvéolos -- paladar blando -- velo del paladar -- dientes --
        labios -- paladar duro -- úvula
\end{consignas}
\end{ejercicio}

\section{Transcripción fonética y descripción articulatoria de los sonidos}

Como pudo observarse en las secciones anteriores, los sistemas de escritura no
representan con exactitud la oralidad. Para graficar la amplia diversidad que
presentan los sonidos del habla y evitar confusiones, se dispone del Alfabeto
Fonético Internacional (AFI), que identifica cada sonido lingüístico humano ---o
\termino{fono}--- en virtud de sus características articulatorias. El AFI, creado
en 1888 y actualizado cada cierta cantidad de años, consiste en un alfabeto
compuesto de símbolos gráficos y diacríticos\footnote{Los diacríticos son signos
gráficos que modifican el valor de un signo de representación fonética.}
aprobados por la Asociación Fonética Internacional para representar los sonidos
del habla de todas las lenguas. Para no confundir su representación, los fonos o
sonidos lingüísticos se escriben entre corchetes; por ejemplo, cuando leemos
\fon{a} o \fon{p} se está hablando de la pronunciación oral y no de los
grafemas.

En una primera aproximación, los sonidos se pueden distinguir en dos grandes
clases: vocales y consonantes. Las \termino{vocales} son aquellos segmentos que
se producen sin obstrucción en el paso del aire, mientras que las
\termino{consonantes} ocurren cuando se produce una obstrucción total o parcial
durante la salida de la corriente de aire.\footnote{Si bien existen sonidos
vocálicos y consonánticos prototípicos, no todos los fonos se ajustan a esta
distinción. No hay consenso para determinar qué parámetro establece la oposición
clara entre una vocal y una consonante.}

\begin{recuadro}[Para ampliar]
El Alfabeto Fonético Internacional, también conocido como IPA (\textit{International
Phonetic Alphabet}) por sus siglas en inglés, es uno de los alfabetos más
utilizados. Sin embargo, no es el único. Para la transcripción del español,
muchos textos académicos han utilizado el alfabeto diseñado por la
\textit{Revista de Filología Española} (RFE). Asimismo, para entornos digitales
se emplea el alfabeto SAMPA (\textit{Speech Assessment Methods Phonetic
Alphabet}), que es compatible con cualquier teclado.\footnote{\textcite{gurlekian2001}
han hecho una adaptación de SAMPA al español.} La necesidad de contar con
alfabetos de este tipo no solo se circunscribe a las áreas de fonética y
fonología. Numerosas disciplinas requieren de instrumentos adecuados para
transcribir los fenómenos orales, entre ellas se encuentran la dialectología, la
enseñanza de lenguas, la etnolingüística y la lingüística computacional.
\end{recuadro}

\subsection{Vocales}

Dado que no hay interrupción en el paso del aire, los segmentos vocálicos se
articulan en relación con tres parámetros:

\begin{itemize}
  \item grado de abertura de la cavidad bucal,
  \item posición de la lengua respecto del paladar,
  \item redondeamiento de los labios.
\end{itemize}

El grado de abertura bucal permite distinguir entre vocales abiertas, medias y
cerradas (y otros grados intermedios). En el trapecio vocálico que acompaña al
AFI (ver figura~\ref{fig:03-vocales}), se puede observar que la máxima abertura
la tienen las vocales de la base \fon{a, ɶ, ɑ, ɒ}, mientras que en las filas
superiores se encuentran las más cerradas \fon{i, y, ɨ, ʉ, ɯ, u}, que se
articulan cuando la lengua se aproxima al paladar.

\figurapendiente{Vocales (AFI). En el esquema, las vocales que se presentan a la
derecha de otra son redondeadas; así por ejemplo \fon{œ} es redondeada respecto
de \fon{ɛ}, y \fon{o} respecto de \fon{ɤ}. Enlace al AFI:
\url{http://www.internationalphoneticalphabet.org}}{fig:03-vocales}[55mm]

Si miramos el trapecio de izquierda a derecha, las columnas muestran la
distinción entre vocales anteriores, centrales y posteriores. Esta clasificación
está determinada por la posición de la lengua. En la producción de las vocales
anteriores, el dorso de la lengua se eleva hacia el paladar (está adelantado),
como ocurre en \fon{i}, mientras que en las posteriores, la lengua se eleva
sobre el velo del paladar (se posiciona hacia atrás), como cuando pronunciamos
\fon{u}. Por último, el redondeamiento de los labios o labialización distingue
entre los pares izquierda y derecha de cada punto del trapecio. La vocal \fon{i}
se articula sin redondeamiento de los labios, a la vez que \fon{y} es producida
con redondeamiento labial. La figura~\ref{fig:03-resonancia} ilustra los
parámetros articulatorios que distinguen las vocales \fon{a, e, i, o, u}.

\figurapendiente{Articulación de vocales. En las imágenes, tomadas a partir de
resonancias magnéticas funcionales, se puede ver la posición de la lengua, la
abertura de la boca y el redondeamiento de los labios en la producción de las
vocales \fon{a, e, i, o, u}. Video completo:
\url{https://www.youtube.com/watch?v=jaIquq_4560} (Matt Edwards,
EdwardsVoice.com)}{fig:03-resonancia}

\begin{ejercicio}\label{ej:03-vocales}
Las siguientes palabras tienen marcado un grafema que representa un segmento
vocálico. Seleccione la descripción articulatoria que lo describe, como se
muestra en el ejemplo resuelto~1.

\begin{tabular}{@{}l l @{\qquad} l@{}}
  1. \forma{v\textbf{a}}      & \fon{a} & a. cerrado anterior no redondeado \\
  2. \forma{raz\textbf{ó}n}   & \fon{o} & b. cerrado posterior redondeado \\
  3. \forma{r\textbf{u}ta}    & \fon{u} & c. abierto central no redondeado \\
  4. \forma{s\textbf{i}n}     & \fon{i} & d. medio cerrado posterior redondeado \\
  5. \forma{\textbf{e}l}      & \fon{e} & e. medio cerrado anterior no redondeado \\
\end{tabular}

\begin{resuelto}
1. \forma{va} \fon{a}: c. abierto central no redondeado.
\end{resuelto}
\end{ejercicio}

\subsection{Consonantes}

Las \termino{consonantes} se producen cuando los articuladores crean alguna
obstrucción durante el paso del aire. Los sonidos consonánticos se clasifican a
partir de cuatro parámetros articulatorios:

\begin{itemize}
  \item modo de articulación,
  \item punto o zona de articulación,
  \item acción de las cuerdas vocales,
  \item acción del velo del paladar.
\end{itemize}

El \termino{modo de articulación} es la manera en la cual sale el aire durante
la producción de la consonante. Esto depende del grado de acercamiento de los
órganos articuladores y del tipo de salida de la corriente de aire. Los órganos
pueden obstruirse completamente, lo que da lugar a sonidos oclusivos, como
\fon{p} en \forma{\textbf{p}az} o \fon{g} en el sonido inicial de
\forma{\textbf{g}ato}. También pueden estrecharse entre sí o aproximarse; así se
generan los sonidos fricativos, como por ejemplo \fon{s} en
\forma{\textbf{s}apo}, y también los aproximantes, como \fon{j} en
\forma{rub\textbf{i}o}. Los sonidos africados incluyen dos fases de articulación
contiguas: una oclusiva y otra fricativa. La única africada del español es la
que se encuentra en \forma{o\textbf{ch}o}. Esta consonante tiene una fase
oclusiva \fon{t} y otra fricativa \fon{ʃ}; por lo tanto, se representa
\fon{t̠͡ʃ}.\footnote{Las consonantes africadas, dado que se producen como
articulaciones dobles, se representan con los símbolos correspondientes unidos
por un diacrítico superior \fon{t̠͡ʃ}.} Si se produce un contacto rápido de los
órganos, se generan los sonidos vibrantes, que pueden ser múltiples (con muchos
contactos de la lengua contra el paladar, como \fon{r} en \forma{pa\textbf{rr}a})
o simples (un solo contacto, como \fon{ɾ} de \forma{pa\textbf{r}a}). Por último,
en los sonidos laterales la lengua genera una constricción central que ocasiona
que el aire fluya por los costados de la cavidad bucal, como ocurre en el sonido
inicial de \forma{\textbf{l}interna} \fon{l}. Los modos de articulación aparecen
representados en las filas del cuadro del AFI
(figura~\ref{fig:03-consonantes}).

El \termino{punto o zona de articulación} sirve para clasificar los sonidos
consonánticos teniendo en cuenta el lugar o la zona en la cual los articuladores
que intervienen en su producción se conectan o aproximan. El cuadro del AFI
(figura~\ref{fig:03-consonantes}) presenta horizontalmente, de izquierda a
derecha, las posibilidades de articulación del sonido desde la parte más
externa de la boca, los labios, hasta el último punto de contacto, que es la
glotis. De esta manera, los sonidos se clasifican en bilabiales (\fon{b} en
inicio de palabra como \forma{\textbf{b}arco}), labiodentales (\fon{f} en
\forma{je\textbf{f}e}), dentales (\fon{θ} en \forma{ca\textbf{z}a} para las
variedades peninsulares no seseantes), alveolares (\fon{t} en
\forma{\textbf{t}alco}), postalveolares (\fon{ʃ} en \forma{llegar} y \forma{yo}
para la variedad rioplatense), palatales (\fon{ɲ} en \forma{a\textbf{ñ}o}) y
velares (\fon{k} en \forma{\textbf{c}abeza}). A pesar de que las articulaciones
de consonantes posteriores (de la zona uvular, faríngea y glotal) no son
frecuentes en español, son comunes en otras lenguas. Los puntos de articulación
aparecen representados en las columnas del cuadro del AFI.

\figurapendiente{Consonantes pulmonares (AFI, 2020). Las casillas en gris
representan articulaciones consideradas imposibles. Las casillas en blanco son
posibles, pero no fueron documentadas hasta el momento}{fig:03-consonantes}[60mm]

\begin{recuadro}[Para ampliar]
Si bien en las lenguas de tipo indoeuropeo (v. \Capref{1}, nota~\ref{fn:01-indoeuropea}),
casi la totalidad de los sonidos son pulmonares, otras lenguas poseen sonidos
con otros mecanismos de iniciación de la corriente de aire. Existen sonidos no
pulmonares que se generan en la glotis. Se dividen en \termino{eyectivos},
cuando el aire sale de la glotis hacia la boca, e \termino{implosivos}, cuando
el aire ingresa. Asimismo, los llamados \textit{clicks}, que presentan algunas
lenguas originarias africanas, se producen por mecanismos de corriente de aire
que se inician en el velo del paladar. En este caso, el aire ingresa y un cierre
bucal genera el \textit{click}. La figura~\ref{fig:03-nopulmonares} incluye la
transcripción de sonidos no pulmonares del AFI.
\end{recuadro}

\figurapendiente{Consonantes no pulmonares (AFI)\protect\pendiente{El original
dice «AFI, 2026»}}{fig:03-nopulmonares}

La \termino{acción de las cuerdas vocales} es un parámetro que distingue los
sonidos en sordos o sonoros. Si se observa el AFI, se puede apreciar que ciertos
puntos y modos de articulación presentan dos alternativas. Los símbolos de la
izquierda representan los sonidos sordos, que son aquellos que se producen sin
actividad de las cuerdas vocales. Los sonidos de la derecha son sonoros, es decir
que se realizan con vibración de las cuerdas vocales. Por ejemplo, el par de
consonantes \fon{b} y \fon{p} presentan el mismo punto y modo de articulación,
pero se distinguen por la acción de las cuerdas: \fon{b} es sonora y \fon{p} es
sorda. Este rasgo permite distinguir palabras como \forma{bala} y \forma{pala}.

Por último, la \termino{acción del velo del paladar} hace referencia a la salida
principal del aire. Si el velo del paladar desciende durante la articulación, la
salida del aire se produce por la nariz, lo que da lugar a sonidos nasales, como
\fon{n} en \forma{\textbf{n}e\textbf{n}e} y \fon{m} en
\forma{\textbf{m}a\textbf{m}á}. Por otra parte, el resto de los sonidos, que se
conocen como orales, se producen cuando el velo asciende y la salida del aire
ocurre por la boca.

En resumen, cada consonante se puede describir a partir de los cuatro
parámetros mencionados. Por ejemplo, el fono \fon{p} es oclusivo (modo de
articulación), bilabial (zona de articulación), sordo (acción de las cuerdas
vocales) y oral (acción del velo del paladar). Se denomina \termino{descripción
articulatoria} a la enumeración de rasgos que configuran un sonido.

\begin{ejercicio}\label{ej:03-transcripcion}
Una cada palabra con su correspondiente transcripción fonética. (Se muestra el
primer ejemplo resuelto).

\begin{tabular}{@{}l @{\qquad} l@{}}
  A. ruta       & 1. \fon{ˈombɾe} \\
  B. por        & 2. \fon{ˈbjento} \\
  C. suerte     & 3. \fon{ba} \\
  D. va         & 4. \fon{ˈsweɾte} \\
  E. Manu Chao  & 5. \fon{manuˈtʃao} \\
  F. viento     & 6. \fon{ˈruta} \\
  G. el         & 7. \fon{poɾ} \\
  H. hombre     & 8. \fon{el} \\
  I. carretera  & 9. \fon{kareˈteɾa} \\
\end{tabular}

\begin{resuelto}
A. ruta -- 6. \fon{ˈruta}.
\end{resuelto}
\end{ejercicio}

\begin{ejercicio}\label{ej:03-articulacion}
\begin{consignas}
  \item Encuentre en la canción de Manu Chao al menos dos palabras que comiencen
        con sonidos cuyo modo de articulación sea oclusivo o fricativo y complete
        el cuadro. (Se muestra el primer ejemplo resuelto).

        \begin{tabular}{@{}Y{4cm} Y{4cm}@{}}
          \toprule
          \textbf{Oclusivos} & \textbf{Fricativos} \\
          \midrule
          \textit{Ej.}: \forma{carretera} & \textit{Ej.}: \forma{se} \\
          & \\
          \bottomrule
        \end{tabular}
  \item Distribuya las siguientes palabras en la tabla de abajo según el punto
        de articulación de la consonante que las inicia.

        \fon{ˈsweɾte} \fon{ˈmanu} \fon{ˈkwando} \fon{ˈbjento} \fon{tʃao}
        \fon{fɾonˈteɾa} \fon{raˈson} \fon{poɾ}

        \begin{tabular}{@{}*{5}{Y{1.9cm}}@{}}
          \toprule
          \textbf{bilabial} & \textbf{labiodental} & \textbf{alveolar} & \textbf{palatal} & \textbf{velar} \\
          \midrule
          \fon{ˈmanu} & & & & \\
          \bottomrule
        \end{tabular}
  \item En los ejercicios a) y b) encontramos palabras cuyo segmento inicial se
        caracteriza por ser bilabial y oclusivo, ¿qué parámetro articulatorio
        permite distinguir estos sonidos?
  \item Encuentre palabras en la canción de Manu Chao que contengan segmentos
        con las siguientes características articulatorias:
        \begin{exe}
          \ex Alveolar oclusiva nasal sonora: \respuesta{\fon{n} en \fon{ˈbjene}}
          \ex Bilabial oclusiva oral sonora:
          \ex Alveolar vibrante múltiple oral sonora:
          \ex Labiodental fricativa oral sorda:
          \ex Alveolar lateral oral sonora:
          \ex Alveolar oclusiva oral sorda:
        \end{exe}
  \item Realice la descripción articulatoria de los siguientes fonos:
        \begin{exe}
          \ex \fon{β}: \respuesta{bilabial fricativa (aproximante) oral
              sonora}\footnote{V. §\ref{sec:03-alofonos} para la aclaración
              sobre la caracterización de este sonido y la diferencia entre
              fricativas y aproximantes.}
          \ex \fon{ɛ}
          \ex \fon{h}
          \ex \fon{x}
          \ex \fon{ɡ}
          \ex \fon{ʒ}
          \ex \fon{ɔ}
        \end{exe}
  \item Transcriba ortográficamente las siguientes palabras:

        \fon{βeˈɾano} \fon{ˈsizmo} \fon{akoˈplaβa} \fon{ˈoŋɡo} \fon{ˈleʧe}
        \fon{ˈsjero} \fon{arkeˈoloɣa} \fon{ˈweso} \fon{ˈpeɾa} \fon{ˈkaɲa}
        \fon{ˈtaksi}

        \begin{resuelto}
        1. \fon{βeˈɾano}: \forma{verano}.
        \end{resuelto}
  \item Transcriba fonéticamente las siguientes palabras:

        \forma{manzana, caño, llave, dura, neurona, pienso, palo, hora, cereza,
        puerro, cama, ley, terapeuta, jota}

        \begin{resuelto}
        1. \forma{manzana}: \fon{manˈsana}.
        \end{resuelto}
\end{consignas}
\end{ejercicio}

\subsection{Los diptongos}

Al observar la transcripción de las palabras \forma{viento} y \forma{suerte}
(\fon{ˈbjento} \fon{ˈsweɾte}), se puede advertir que la primera vocal del
diptongo se representa de manera diferente de la vocal cuando está sola en la
sílaba (este concepto se desarrolla más adelante en §\ref{sec:03-silaba}). Si
dentro de una sílaba se incluye un diptongo, el núcleo silábico es complejo:
una de las vocales de la sílaba es plena, mientras que las que la preceden o
anteceden tienden a ser aproximantes o transicionales, por lo que su
articulación se acerca, sin llegar a serlo, a la de una consonante. Son
\termino{semiconsonantes} los sonidos \fon{j} y \fon{w} que están en primera
posición del diptongo, como en las palabras \forma{pie} y \forma{suela}, o en
posición final, como en \forma{peine} y \forma{deuda}.

\subsection{La representación del acento}

Se puede observar que en la transcripción de las palabras aparece el diacrítico
\fon{ˈ} que antecede a un segmento o sílaba de la palabra. Se trata de la
transcripción del \termino{acento léxico}. El acento es el realce que se realiza
sobre una sílaba de la palabra en contraste con otras. Debe distinguirse el
\termino{acento prosódico}, que es propio de la oralidad, del acento
ortográfico o \termino{tilde}. En el nivel prosódico, todas las palabras tienen
acento;\footnote{En §\ref{sec:03-suprasegmentales} se explican algunas
excepciones a esta afirmación. A su vez, se ampliarán los detalles sobre este
fenómeno y sus implicancias.} en cambio, a nivel ortográfico, solo en ciertas
palabras y siguiendo reglas específicas, se destacan algunas sílabas acentuadas
con un signo (´). Como para la transcripción fonética lo que cuenta es la
pronunciación oral, el acento prosódico se marca de forma obligatoria siempre e
independientemente de la ortografía que rige en cada lengua.

\section{Los sonidos distintivos o fonemas}

Las lenguas humanas utilizan y combinan un número relativamente limitado de
sonidos para la expresión de todas sus unidades: palabras, enunciados y textos.
Algunos de estos sonidos son distintivos y se denominan \termino{fonemas}. Cada
lengua dispone de una cantidad restringida de unidades distintivas o fonemas
que, si bien carecen de significado, permiten distinguir los significados de
unidades mayores. Por ejemplo, en español las palabras \forma{sol} y \forma{sal}
se diferencian entre sí por los segmentos vocálicos \fonol{o} y \fonol{a},
respectivamente. Esto evidencia que para el español estos sonidos son fonemas,
ya que su contraste permite distinguir dos palabras de la lengua. Lo mismo
ocurre con los pares \forma{palo-polo}, \forma{dato-pato} y \forma{lado-lodo}.

Estos pares de palabras conocidos como \termino{pares mínimos} posibilitan
identificar si dos sonidos constituyen o no fonemas de una lengua mediante la
prueba de conmutación en un mismo contexto. Se trata de una prueba diagnóstica
para reconocer el inventario de unidades opositivas de una lengua (fonemas), en
la medida en que la aparición de uno u otro fonema contrasta significados. Las
palabras \forma{toro} y \forma{coro} constituyen un par mínimo,\footnote{V.
\Capref{1}, §\ref{sec:01-paresminimos}.} ya que se diferencian únicamente por la
consonante inicial. La aparición del fonema \fonol{t} o \fonol{k} en posición
inicial conlleva un cambio de significado producido por la aparición del fono
\fon{t} o \fon{k} en el mismo contexto \fon{\_oro}. Con el descubrimiento de este
par mínimo y ante la posibilidad de contrastar significados para las secuencias
\fon{ˈtoɾo} y \fon{ˈkoɾo}, es posible concluir que los sonidos \fon{t} y \fon{k}
son fonemas distintos, ya que dentro del sistema fonológico del español los
fonos \fon{t} y \fon{k} tienen una distribución contrastiva, es decir, que su
presencia en el mismo contexto opone significados diferentes.

Es importante destacar que los fonemas son unidades abstractas: no se trata de
sonidos en sí, sino de entidades mentales que forman parte de la competencia
lingüística de los hablantes. Se forman a partir de una serie de
\termino{rasgos distintivos} que permiten que los reconozcamos e identifiquemos
frente a otros fonemas. La identidad de un fonema se constituye a partir de las
posibilidades opositivas que genera en el sistema frente a otras unidades. Las
pruebas diagnósticas para el descubrimiento del inventario fonológico de una
lengua tienen sentido únicamente dentro del sistema de cada lengua particular, ya
que cada una tiene su propio conjunto de fonemas. La cantidad de fonemas y su
combinatoria dependen del sistema lingüístico que se esté estudiando.

\begin{recuadro}[Para ampliar]
Los fonemas constituyen un número limitado de elementos que están representados
por rasgos que los diferencian unos de otros en la realidad lingüística mental.
Por ejemplo, el fonema \fonol{p} se puede caracterizar como \sig{consonántico},
\sig{oral}, \sig{oclusivo}, \sig{bilabial}, \sig{sordo}. Los cuatro primeros
rasgos los comparte con \fonol{b}, que es también \sig{consonántico},
\sig{oral}, \sig{oclusivo} y \sig{bilabial}, pero del que se diferencia porque
\fonol{b} es \sig{sonoro}; en cambio, respecto de \fonol{m}, los rasgos que
comparte son solo tres: \sig{consonántico}, \sig{oclusivo},
\sig{bilabial}, ya que \fonol{m} es \sig{nasal} y \sig{sonoro}. A pesar de que
cada lengua difiera en la cantidad de fonemas que posee, todas las lenguas
incluyen un grupo pequeño de fonemas que tienen una amplia posibilidad
combinatoria entre sí. En español, la cantidad de fonemas varía según la
variedad lingüística que se estudie. Se cuentan unos 24 fonemas: diecinueve son
consonantes y cinco son vocales. En algunas variedades ese número cambia por el
seseo y el yeísmo, a los que ya nos referimos en §\ref{sec:03-grafemas}.
\end{recuadro}

\begin{ejercicio}\label{ej:03-paresminimos}
\begin{consignas}
  \item Observe el siguiente corpus e identifique un par mínimo.
  \item Proporcione pares mínimos para las otras palabras del corpus.
\end{consignas}

\fon{mas} \fon{ba} \fon{ˈmanu} \fon{ˈbjento} \fon{la} \fon{lo} \fon{sin}
\fon{ˈombɾe} \fon{ˈsweɾte} \fon{poɾ} \fon{ˈambɾe} \fon{se}

\begin{resuelto}[Punto (a) resuelto]
Par mínimo presente en el corpus: \fon{ˈombɾe} y \fon{ˈambɾe}.
\end{resuelto}
\end{ejercicio}

\begin{ejercicio}\label{ej:03-ue}
\begin{consignas}
  \item Observe el siguiente corpus del español. ¿Se puede identificar algún par
        mínimo?
  \item ¿Qué conclusión se puede sacar de la distribución de los fonos \fon{u}
        y \fon{e} en el sistema fonológico del español?
\end{consignas}

\fon{pus} \fon{seɾ} \fon{ˈluʧe} \fon{ˈpeɣo} \fon{saˈβeɾ} \fon{ˈleʧe}
\fon{ˈuniko} \fon{suˈβiɾ} \fon{suɾ} \fon{enˈtjero} \fon{pes}
\end{ejercicio}

\begin{ejercicio}\label{ej:03-aleman}
\begin{consignas}
  \item Observe el siguiente corpus del alemán. ¿Se puede identificar algún par
        mínimo?
  \item ¿Qué conclusión se puede sacar de la distribución de los fonos \fon{uː}
        y \fon{ʊ} en el sistema fonológico del alemán?
\end{consignas}

\begin{tabular}{@{}lll@{}}
  \fon{ˈp͡fuːl} \sig{estanque}    & \fon{ˈʃpuːkən} \sig{embrujado} & \fon{ˈbaʊm} \sig{árbol} \\
  \fon{ˈhuːn} \sig{pollito}       & \fon{ˈʊntən} \sig{abajo}       & \fon{ˈbʊʃ} \sig{arbusto} \\
  \fon{ˈbʊnt} \sig{colorido}      & \fon{ˈbuːx} \sig{libro}        & \fon{ˈhaʊs} \sig{casa} \\
  \fon{ˈʃpʊkən} \sig{escupir}     & & \\
\end{tabular}
\end{ejercicio}

\section{Los alófonos o variantes de un fonema}\label{sec:03-alofonos}

El sistema fonológico de cada lengua en particular cuenta con un grupo reducido
de fonemas cuya combinación genera unidades mayores significativas; así, por
ejemplo, la combinación de los fonemas \fonol{a} \fonol{l} \fonol{a} produce la
palabra \forma{ala}, mientras que los fonemas \fonol{a} \fonol{s} \fonol{a}
forman la palabra \forma{asa}. Sin embargo, dado que, como ya dijimos, el fonema
es una unidad abstracta, la producción real del sonido lingüístico no ocurre
siempre de la misma manera. En otras palabras, la pronunciación de un fonema
puede variar según el contexto en el que aparezca.

Los sonidos de una lengua nunca se pronuncian exactamente igual, incluso si se
trata de un mismo hablante. Existen factores propios del habla que influyen en
la producción de los sonidos lingüísticos: el contexto, la velocidad, la
intención comunicativa, la variación dialectal o sociolectal, entre otros.

Cada lengua puede presentar variaciones sistemáticas en la pronunciación de un
mismo fonema. Este fenómeno es conocido como \termino{alofonía}. Como puede
observarse en el corpus del ejercicio~\ref{ej:03-bd}, en la palabra
\forma{babero} \fon{baˈβeɾo} podemos reconocer dos segmentos distintos \fon{b} y
\fon{β}, que son alófonos del fonema \fonol{b}, porque su mutación no implica un
cambio de significado. Ambos segmentos son bilabiales orales sonoros, pero
difieren en el modo de articulación, ya que \fon{b} es oclusivo y \fon{β}
representa un sonido fricativo o aproximante.\footnote{Si bien se utilizan los
símbolos del AFI que corresponden al modo de articulación fricativo para las
consonantes intervocálicas \fon{β ð ɣ}, en español estos alófonos se realizan
como aproximantes, ya que se producen con un grado un tanto mayor de apertura
que los sonidos fricativos. A pesar de que algunas transcripciones utilizan un
diacrítico para marcar la diferencia entre una consonante fricativa \fon{ð} y
una aproximante \fon{ð̞}, en este capítulo se simplificará la notación y se
tomarán los símbolos \fon{β ð ɣ} como representaciones de sonidos
aproximantes.} Entonces, la distinción entre el rasgo oclusivo o aproximante en
estos dos fonos en particular no es contrastiva para el español e incluso esta
variación, aunque sea sistemática, resulta imperceptible para los hablantes no
entrenados.

\begin{ejercicio}\label{ej:03-bd}
\begin{consignas}
  \item Lea en voz alta las siguientes palabras y observe la pronunciación de
        las consonantes iniciales y medias marcadas en negrita. ¿Son similares o
        es posible percibir una diferencia en su articulación?

        \begin{tabular}{@{}l@{\qquad}l@{}}
          \forma{\textbf{B}a\textbf{b}ylon} & \forma{\textbf{D}eu\textbf{d}a} \\
          \forma{\textbf{B}e\textbf{b}ida}  & \forma{\textbf{D}e\textbf{d}ica\textbf{d}o} \\
          \forma{\textbf{B}a\textbf{b}ero}  & \forma{\textbf{D}e\textbf{d}al} \\
        \end{tabular}
  \item A continuación se presenta la transcripción fonética de las palabras del
        ejercicio anterior, ¿qué se puede observar?

        \fon{ˈbaβilon} \fon{baˈβeɾo} \fon{deˈðal} \fon{ˈdeṷða} \fon{deðiˈkaðo}
        \fon{beˈβiða}
  \item Realice la descripción articulatoria de \fon{b}, \fon{β}, \fon{d} y
        \fon{ð}. (Se muestra el primer ejemplo resuelto).
        \begin{resuelto}[Ejemplos resueltos]
        \fon{b} bilabial oclusiva oral sonora y \fon{β} bilabial aproximante
        (fricativa) oral sonora.
        \end{resuelto}
  \item ¿Cuál es el contexto que explica la aparición de \fon{b} \fon{β} y
        \fon{d} \fon{ð}?
\end{consignas}
\end{ejercicio}

\section{Distribución contrastiva y distribución complementaria}\label{sec:03-distribucion}

El concepto de contexto es clave para comprender cómo se distribuyen los
sonidos de una lengua. En fonética y fonología, cuando se habla de contexto, se
hace referencia a los segmentos (sonidos) o propiedades suprasegmentales (como
el acento o el tono, ver §\ref{sec:03-suprasegmentales}) que rodean al elemento
que se está analizando. Puede tratarse de una vocal o de una consonante que
antecede o sucede a un segmento. Por ejemplo, en la palabra \forma{sabe}, el
fonema \fonol{b} está en contexto intervocálico, ya que aparece entre las
vocales \fon{a} y \fon{e}, lo que hace que, como vimos, su articulación sea
aproximante \fon{β}. Asimismo, el contexto puede referirse a una unidad
suprasegmental, como el acento. Por ejemplo, en el par de palabras
\forma{paso-pasó}, la vocal \fonol{o} en el primer caso se presenta en una
sílaba no acentuada, mientras que en \forma{pasó} la vocal está en una sílaba
tónica o acentuada.

Según ya hemos dicho, los fonemas no son equivalentes a un sonido «real». Se
trata de unidades abstractas que pueden manifestarse concretamente mediante uno,
dos o más alófonos. En algunos casos, la presencia de uno u otro alófono depende
del contexto en el que aparezcan. Por ejemplo, en español \fon{d} y \fon{ð} son
alófonos de un mismo fonema: el alófono \fon{ð} aparece en contexto
intervocálico (entre vocales) y el alófono \fon{d} después de pausa o después de
consonante nasal, como ocurre en \fon{deðiˈkaðo} y \fon{kondoˈminjo}. Esto
significa que en el sistema fonológico del español estos fonos están en
\termino{distribución complementaria}. Juntos forman una unidad abstracta: el
fonema \fonol{d}. La diferencia material que hay entre estos sonidos no genera
oposiciones de significado para el español; ambos alófonos conforman una única
entidad mental. Este fenómeno es recurrente en español y, además de con
\fonol{b} y \fonol{d}, lo podemos observar también en el fonema \fonol{g}, cuyos
alófonos son \fon{g}-\fon{ɣ}, como en \fon{ˈgato} y \fon{aɣoˈtaðo}. Los ejemplos
muestran que, cuando el contexto fonético cambia, el alófono correspondiente se
realiza en función de las vocales adyacentes. Esto se debe a la coarticulación
de las vocales con las consonantes que les siguen: los órganos no llegan a
cerrarse del todo cuando tienen que producir una consonante entre dos vocales;
por lo tanto, no pueden realizarse como oclusivos.

Sin embargo, cada lengua presenta su propia configuración fonológica. En
inglés, por ejemplo, los mismos sonidos \fon{d} y \fon{ð} anteriores no son
alófonos, sino fonemas distintos porque distinguen significados. El par mínimo
que forman las palabras \forma{day} \fon{deɪ} \sig{día} y \forma{they}
\fon{ðeɪ} \sig{ellos} muestra que en el mismo contexto la producción de un fono
o de otro genera un cambio de significado. En el sistema fonológico del inglés
\fonol{d} y \fonol{ð} son dos fonemas distintos; por lo tanto, los sonidos o
fonos \fon{d} y \fon{ð} se encuentran en \termino{distribución contrastiva},
como ocurre con los fonos \fon{e} y \fon{u} en español. En la
figura~\ref{fig:03-espectrograma} se muestra la imagen de un espectrograma y un
oscilograma en el que aparece una variación alofónica.

\figurapendiente{Alófonos del fonema \fonol{b}. Espectrograma y oscilograma de
la palabra \textit{bebé}. Se puede observar que la producción de la primera
consonante no es similar a la segunda, que está en posición intervocálica. Se
trata de una consonante oclusiva y otra aproximante,
respectivamente}{fig:03-espectrograma}

\begin{ejercicio}\label{ej:03-sz-espanol}
\begin{consignas}
  \item Observe el siguiente corpus del español. Defina las características
        articulatorias de \fon{s} y \fon{z}.
  \item ¿El corpus presenta algún par mínimo?
  \item ¿Qué distribución tienen los fonos \fon{s} y \fon{z} en español?
        Justifique y explicite los contextos en los que aparecen.
\end{consignas}

\begin{tabular}{@{}llll@{}}
  \fon{dezˈðen} & \fon{sin}         & \fon{sin ˈθeɾo} & \fon{oˈsaða} \\
  \fon{raˈson}  & \fon{konsuˈmizmo} & \fon{aβizˈmal}  & \fon{fizɣoneˈar} \\
\end{tabular}

\begin{resuelto}[Punto (a) resuelto]
El sonido \fon{s} es alveolar, fricativo, oral, sordo y \fon{z} es alveolar,
fricativo, sonoro.
\end{resuelto}
\end{ejercicio}

\begin{ejercicio}\label{ej:03-sz-frances}
\begin{consignas}
  \item Observe el siguiente corpus del francés. ¿Presenta algún par mínimo?
  \item ¿Qué distribución tienen los fonos \fon{s} y \fon{z} en francés?
        Justifique y explicite los contextos en los que aparecen.
\end{consignas}

\begin{tabular}{@{}llll@{}}
  \fon{pwasɔ̃} \sig{pez}     & \fon{dəsu} \sig{debajo}   & \fon{dezɛʀ} \sig{desierto} & \fon{sak} \sig{bolsa} \\
  \fon{gɑzɔ̃} \sig{césped}   & \fon{fos} \sig{fosa}      & \fon{pwazɔ̃} \sig{veneno}   & \fon{zon} \sig{zona} \\
  \fon{desɛʀ} \sig{postre}  & \fon{pusjɛʀ} \sig{polvo}  & \fon{vizɔ̃} \sig{visón}     & \\
\end{tabular}

\begin{resuelto}[Punto (a) resuelto]
Pares mínimos del corpus: \fon{pwasɔ̃} \sig{pez} / \fon{pwazɔ̃} \sig{veneno};
\fon{dezɛʀ} \sig{desierto} / \fon{desɛʀ} \sig{postre}.
\end{resuelto}
\end{ejercicio}

\begin{ejercicio}\label{ej:03-quechua}
\begin{consignas}
  \item Observe los siguientes corpus del quechua y del español. Caracterice las
        vocales \fon{i}, \fon{e}, \fon{o} y \fon{u} desde el punto de vista
        articulatorio.
  \item Indique el tipo de distribución y los contextos de aparición de las
        vocales \fon{i} y \fon{e} para el quechua.
  \item Indique el tipo de distribución y los contextos de aparición de las
        vocales \fon{i} y \fon{e} para el español.
  \item Indique el tipo de distribución y los contextos de aparición de las
        vocales \fon{o} y \fon{u} para el quechua.
  \item Indique el tipo de distribución y los contextos de aparición de las
        vocales \fon{o} y \fon{u} para el español.
\end{consignas}

\textit{Quechua}

\begin{tabular}{@{}llll@{}}
  \fon{weqe} \sig{lágrima}   & \fon{suti} \sig{nombre}      & \fon{qespi} \sig{vidrio}    & \fon{niy} \sig{decir} \\
  \fon{apiy} \sig{agarrar}   & \fon{pisiŋga} \sig{maíz}     & \fon{qožqe} \sig{dinero}    & \fon{sapi} \sig{raíz} \\
  \fon{puriy} \sig{andar}    & \fon{t͡ʃeqɑ} \sig{cierto}     & \fon{urpila} \sig{paloma}   & \fon{šiši} \sig{hormiga} \\
  \fon{t͡ʃoqɑ} \sig{arrojar}  & \fon{punčaw} \sig{día}       & \fon{rumi} \sig{piedra}     & \fon{ušpa} \sig{cenizas} \\
  \fon{žoqe} \sig{izquierda} & \fon{aku} \sig{harina}       & \fon{ñuñu} \sig{seno}       & \fon{ažqo} \sig{perro} \\
\end{tabular}

\medskip
\textit{Español}

\begin{tabular}{@{}llll@{}}
  \fon{ˈmisa}       & \fon{seˈkreto}     & \fon{ˈmesa}      & \fon{termiˈno} \\
  \fon{ˈpesa}       & \fon{miˈlimetro}   & \fon{ˈpisa}      & \fon{ˈkeso} \\
  \fon{nesesiˈðad}  & \fon{senˈtiðo}     & \fon{puˈðjente}  & \fon{ˈmuðo} \\
  \fon{xiˈrafa}     & \fon{konoˈsjendo}  & \fon{ˈoxos}      & \fon{ˈkulpa} \\
  \fon{guˈsano}     & \fon{luˈɣar}       & \fon{ˈpuso}      & \fon{kurjosiˈðad} \\
  \fon{ˈposo}       & \fon{lusjeɾˈnaɣa}  & \fon{multiˈtud}  & \fon{ˈmoðo} \\
\end{tabular}

\begin{resuelto}[Ejemplo resuelto del punto (a)]
La vocal \fon{i} es anterior, cerrada y no redondeada.
\end{resuelto}
\end{ejercicio}

\begin{ejercicio}\label{ej:03-ruso}
\begin{consignas}
  \item Observe el siguiente corpus del ruso. ¿Se puede identificar algún par
        mínimo dentro del corpus?
  \item ¿Qué parámetro articulatorio diferencia \fon{ɨ} de \fon{i}?
  \item ¿Qué distribución tienen estos fonos en ruso?
  \item ¿Cuál es el contexto de aparición de cada sonido?
\end{consignas}

\begin{tabular}{@{}llll@{}}
  \fon{dˠɨm} \sig{humo}         & \fon{pastuxʲi} \sig{pastores} & \fon{tenʲisʲist} \sig{tenistas} & \fon{mʲixka} \sig{suave} \\
  \fon{pˠɨtkə} \sig{tortura}    & \fon{stolˠɨ} \sig{mesas}      & \fon{zakonˠɨ} \sig{leyes}       & \fon{vʲitɐmʲin} \sig{vitamina} \\
  \fon{pˠɨˈtatsə} \sig{tratar}  & \fon{ʌlfʌvʲit} \sig{alfabeto} & \fon{zubˠɨ} \sig{dientes}       & \fon{kʲijev} \sig{Kiev} \\
  \fon{vʲit} \sig{monedas}      & \fon{nosˠɨ} \sig{narices}     & & \\
\end{tabular}
\end{ejercicio}

\section{Variación libre}

La distribución que presentan dos alófonos de un mismo fonema es complementaria
si su producción puede predecirse de acuerdo con la posición silábica, el
contacto con otros sonidos o el contexto de aparición, entre otros factores. En
estos casos, la aparición de uno u otro es excluyente y no son intercambiables:
el contexto determina si ocurre uno u otro alófono.

No obstante, existen variaciones alofónicas que no dependen del contexto
fonético, sino de las características de los hablantes. Un mismo hablante puede
pronunciar de una manera un sonido cuando su estilo es más cuidado o formal y de
otra cuando está en una situación más informal. Por ejemplo, en el español de la
Argentina, la realización de \fonol{s} en final de palabra puede alternar entre
tres posibles realizaciones: alveolar, aspirada o directamente elidida. El
sintagma \forma{los amigos} puede pronunciarse entonces \fon{losaˈmiɣos},
\fon{lohaˈmiɣoh}, \fon{loaˈmiɣo}, respectivamente.

También un mismo fonema puede pronunciarse de distinta manera según la región
geográfica a la que pertenece un hablante. Por ejemplo, la realización de la
consonante rótica \fonol{r} (la que se pronuncia en \forma{co\textbf{rr}o}) en
Argentina es muy sensible a la variación geográfica. En muchas regiones del
país, principalmente del noroeste, la articulación de esta consonante es de tipo
asibilada \fon{ř},\footnote{Un sonido asibilado se aproxima a un silbido. No
existe un símbolo exacto para la descripción de este sonido en el AFI. En la
bibliografía, se lo encuentra como \fon{ř}.} mientras que en la región de
Buenos Aires, parte de Santa Fe y Patagonia se articula como vibrante múltiple
\fon{r}. De igual forma, hay zonas en las que conviven ambas pronunciaciones y
se distribuyen según las condiciones sociológicas de los hablantes, como se
muestra en el ejercicio~\ref{ej:03-video}.

La coexistencia de estas alternancias en la pronunciación de ciertos fonemas se
puede explicar a partir de factores de índole sociolectal (variaciones sociales
de los hablantes dentro de una misma comunidad, como el sexo, el nivel de
escolarización, la edad, entre otros), de registro (contextos formales vs.
informales) y dialectales (que dependen de la región
geográfica).\footnote{Estas diferencias de lectos se introdujeron en el
\capref{2}.} A este tipo de variación alofónica se la conoce como
\termino{variación libre} o \termino{estilística}.

\begin{ejercicio}\label{ej:03-bsas}
\begin{consignas}
  \item Observe el siguiente corpus del español de Buenos Aires. ¿Qué rasgos
        fonéticos caracterizan a los sonidos \fon{ʃ} y \fon{ʒ}? Realice su
        descripción articulatoria.
  \item ¿Hay pares mínimos en el corpus? ¿En qué contexto aparecen estos
        segmentos?
  \item ¿Cuál es la distribución de estos sonidos en el español de Buenos Aires?
  \item Escuche los audios e identifique a qué hablante corresponde cada corpus.
\end{consignas}

\imagenpendiente[12mm]{Audios A y B del ejercicio}

\begin{tabular}{@{}l Y{2.4cm} Y{2.4cm}@{}}
  \toprule
  Transcripción ortográfica & \multicolumn{2}{c}{Español de Buenos Aires} \\
  \cmidrule(l){2-3}
  & Hablante 1 & Hablante 2 \\
  \midrule
  Río        & \fon{ˈrio}     & \fon{ˈrio} \\
  Desmayo    & \fon{desˈmaʃo} & \fon{desˈmaʒo} \\
  Carretera  & \fon{kareˈteɾa} & \fon{kareˈteɾa} \\
  Llave      & \fon{ˈʃaβe}    & \fon{ˈʒaβe} \\
  Boya       & \fon{ˈboʃa}    & \fon{ˈboʒa} \\
  Callo      & \fon{ˈkaʃo}    & \fon{ˈkaʒo} \\
  Pollo      & \fon{ˈpoʃo}    & \fon{ˈpoʒo} \\
  Parrilla   & \fon{paˈriʃa}  & \fon{paˈriʒa} \\
  \bottomrule
\end{tabular}

\begin{resuelto}[Punto (a) resuelto]
El fono \fon{ʃ} es postalveolar, fricativo, oral, sordo y el fono \fon{ʒ} es
postalveolar, fricativo, oral, sonoro.
\end{resuelto}
\end{ejercicio}

\begin{ejercicio}\label{ej:03-catamarca}
\begin{consignas}
  \item Observe el siguiente corpus del español de
        Catamarca.\footnote{Agradezco a la Prof. Alejandra Seco por su
        colaboración para elaborar este corpus.} ¿Cuál es la distribución de los
        sonidos \fon{ř} y \fon{r} en el español de Catamarca?
  \item ¿Qué motiva la aparición de los fonos en cada muestra?
  \item Escuche los audios e identifique a qué hablante corresponde cada corpus.
\end{consignas}

\imagenpendiente[12mm]{Audios A y B del ejercicio}

\begin{tabular}{@{}l Y{2.4cm} Y{2.4cm}@{}}
  \toprule
  Transcripción ortográfica & \multicolumn{2}{c}{Español de Catamarca} \\
  \cmidrule(l){2-3}
  & Hablante 1 & Hablante 2 \\
  \midrule
  Río        & \fon{ˈřio}      & \fon{ˈrio} \\
  Desmayo    & \fon{desˈmajo}  & \fon{desˈmajo} \\
  Carretera  & \fon{kařeˈteɾa} & \fon{kareˈteɾa} \\
  Llave      & \fon{ˈjaβe}     & \fon{ˈjaβe} \\
  Pollo      & \fon{ˈpojo}     & \fon{ˈpojo} \\
  Parrilla   & \fon{paˈřija}   & \fon{paˈrija} \\
  Borra      & \fon{ˈbořa}     & \fon{ˈbora} \\
  Carro      & \fon{ˈkařo}     & \fon{ˈkaro} \\
  \bottomrule
\end{tabular}
\end{ejercicio}

\begin{ejercicio}\label{ej:03-video}
\begin{consignas}
  \item Lea la siguiente transcripción fonética. ¿Qué ocurre con los fonos
        \fon{r} y \fon{ř}?
  \item ¿Qué puede motivar la aparición de cada sonido?
  \item Realice la transcripción ortográfica del fragmento.
  \item Escuche el video para analizar el fenómeno.\pendiente{Falta el enlace
        al video}
\end{consignas}

\fon{ˈes to es un ˈři o ˈes te es el ka ˈnal de la ma ˈrja no mo ˈre no ˈtje ne
kas ˈka ða a ˈo ra ˈmi ren lo ke ez la ma ˈrja no mo ˈre no las ˈka jes ke ˈβje
nen de a ˈri βa ðe la ma ˈrja no mo ˈre no la ˈpra ðo eh im po ˈsi βle tran si
ˈtar la ˈmi ren lo ke eh ˈes ta ez la ˈro xas ˈmi ren lo ke ez la ˈro xas}

\begin{resuelto}[Punto (a) resuelto]
Los alófonos \fon{ř} y \fon{r} son producidos por el mismo hablante en
distintas partes de su enunciado.
\end{resuelto}
\end{ejercicio}

\begin{ejercicio}\label{ej:03-yendo}
Observe el error ortográfico presente en la siguiente imagen. ¿Con qué
característica fonológica dialectal se lo puede vincular?

\imagenpendiente{Imagen del ejercicio (error ortográfico con \textit{yendo})}
\end{ejercicio}

\section{Procesos fonológicos}\label{sec:03-procesos}

Cuando hablamos, los sonidos se producen dentro de secuencias cada vez más
complejas. De esta manera, la cadena de habla está formada por sonidos que
presentan características combinatorias cuando se articulan conjuntamente, es
decir, cuando se coarticulan. No es lo mismo estudiar cada sonido de forma
aislada que cuando se produce en un contexto particular. Existen sonidos que al
combinarse con otros se modifican como resultado de esa coarticulación.

Se denomina \termino{proceso fonológico} a toda modificación sistemática que se
produce en un fonema cuando ocurre en determinado contexto. La fonética
combinatoria se ocupa de estudiar estos fenómenos. Si bien existen numerosos
procesos fonológicos, en este apartado nos ocuparemos de dos de los más
frecuentes en español: la neutralización y la asimilación.

Hualde (2013, p.~93)\pendiente{Hualde 2013 no está en la bibliografía (¿2014?)}
define \termino{neutralización} como «situaciones en que dos sonidos que en
general se encuentran en oposición fonémica dejan de contrastar en un contexto
específico». Para ampliar este concepto, vamos a partir de lo que se explicó
anteriormente. Para poder establecer que dos sonidos son contrastivos en una
lengua, se debe encontrar un par mínimo que lo verifique. Si la conmutación de
dos segmentos en la misma posición del mismo contexto resulta en un cambio de
significado, se trata de dos sonidos con distribución contrastiva. Este sería el
caso de \fon{b} y \fon{p}; gracias a la existencia de pares mínimos como
\fon{ˈbata} y \fon{ˈpata} o \fon{ˈpeso} y \fon{ˈbeso} se puede afirmar que se
trata de dos fonemas distintos del español: \fonol{b} y \fonol{p}.

Sin embargo, dos sonidos que contrastan en un contexto fonético pueden no
hacerlo en otro. Así, los mismos dos fonos anteriores, según veremos, en otro
contexto no contrastan. Por ejemplo, en la palabra \forma{óptica} la realización
del segundo segmento puede variar sin alterar el significado: \fon{ˈobtika}
\fon{ˈoptika}. Se observa que en la posición de final de sílaba, conocida como
coda silábica, estos dos fonos pierden su valor opositivo. La representación de
este proceso de neutralización se realiza a través de un
\termino{archifonema}, que consiste en una representación en letra mayúscula
del fonema, y que indica el rasgo en común que tienen los segmentos que se
neutralizan; en el ejemplo, la representación fonológica sería
\fonol{ˈoBtika}.

La \termino{asimilación} es un fenómeno articulatorio que ocurre cuando un
sonido incorpora rasgos de otro sonido del contexto cercano. El ejemplo más
claro de esto es lo que ocurre con las nasales en español en contexto
preconsonántico. Si observamos los ejemplos de \exref{ex:03-nasales}, podemos
concluir que la nasal toma el punto de articulación de la consonante que le
sigue. En el caso de \exref{ex:03-nasales-a}, la nasal se produce como bilabial,
dada la adyacencia con \fon{p}; en \exref{ex:03-nasales-b}, se articula como
labiodental por el contacto con \fon{f}. Los casos de \exref{ex:03-nasales-c} y
\exref{ex:03-nasales-d} ilustran el mismo fenómeno: la nasal se articula como
dental/alveolar frente a \fon{d} y como velar antes de \fon{g}.

\begin{exe}
  \ex\label{ex:03-nasales}
    \begin{xlist}
      \ex\label{ex:03-nasales-a} \forma{un poco}  \fon{umˈpoko}
      \ex\label{ex:03-nasales-b} \forma{un faro}  \fon{uɱˈfaɾo}
      \ex\label{ex:03-nasales-c} \forma{un dólar} \fon{unˈdolaɾ}
      \ex\label{ex:03-nasales-d} \forma{un gato}  \fon{uŋˈgato}
    \end{xlist}
\end{exe}

Este tipo de asimilación es la más frecuente en las lenguas; se trata de la
asimilación regresiva, ya que los rasgos de un segmento son asimilados en el
segmento anterior. Cuando la asimilación es progresiva, la direccionalidad
cambia y es el primer segmento el que afecta al siguiente. Según
\textcite{ullman2006}, el fonema nasal también se ve involucrado en procesos de
asimilación progresiva. Este tipo de proceso se da en palabras que incluyen la
secuencia de fonemas \fonol{d} + \fonol{n} como en \forma{llamadnos}, o a través
de fronteras léxicas como en \forma{red nacional}. Según el autor, en muchas
comunidades hispanohablantes el fonema nasal se realiza como apicodental
\fon{n̪} debido a la influencia de \fon{d}, como sucede en
\fon{redn̪asjoˈnal}.\pendiente{En (1) el original transcribe [umpo'ko],
[uɱ'faro], [undo'lar]; se corrigieron el acento y las vibrantes}

\begin{ejercicio}\label{ej:03-griego}
\begin{consignas}
  \item Observe los siguientes corpus del español y del griego. ¿Qué rasgo
        fonético permite distinguir las consonantes \fon{ɾ} y \fon{r}?
  \item ¿Qué distribución presentan estos sonidos en español? ¿Hay un par
        mínimo?
  \item ¿Ocurre esto en todos los contextos?
  \item ¿Cómo se distribuyen estos sonidos en griego? Justifique a través de los
        contextos de aparición.
  \item Formule una conclusión sobre la distribución de estos fonos en ambas
        lenguas.
  \item Represente fonológicamente las palabras \forma{acusarte} y
        \forma{dador}.
\end{consignas}

\textit{Español}

\begin{tabular}{@{}llll@{}}
  \fon{ˈpero}  & \fon{poˈɾoto} & \fon{akuˈsarte} & \fon{ˈpaɾa} \\
  \fon{ˈroxo}  & \fon{daˈðoɾ}  & \fon{ˈpeɾo}     & \fon{akuˈsaɾte} \\
  \fon{ˈpweðo} & \fon{daˈðor}  & \fon{ˈraβja}    & \fon{ˈpwero} \\
  \fon{ˈpara}  & & & \\
\end{tabular}

\medskip
\textit{Griego}

\begin{tabular}{@{}llll@{}}
  \fon{ˈartos} \sig{pan}           & \fon{kaˈɾoto} \sig{zanahoria} & \fon{ˈɾafi} \sig{armario} & \fon{arˈluba} \sig{tontería} \\
  \fon{ˈorkos} \sig{juramento}     & \fon{ˈzaɾi} \sig{dado}        & \fon{ˈturta} \sig{tarta}  & \fon{aˈɾiða} \sig{pie} \\
  \fon{maˈɾuli} \sig{repollo}      & \fon{ˈkerveɾos} \sig{cerbero} & & \\
\end{tabular}

\begin{resuelto}[Punto (a) resuelto]
Las consonantes \fon{ɾ} y \fon{r} son alveolares, orales y sonoras y se
diferencian por su modo de articulación: \fon{ɾ} es vibrante simple y \fon{r}
es vibrante múltiple.
\end{resuelto}
\end{ejercicio}

\begin{ejercicio}\label{ej:03-arabe}
\begin{consignas}
  \item Observe el siguiente corpus de árabe estándar. En cada cuadro se
        encuentra la transcripción fonética y su traducción. Las palabras están
        separadas con un guion. A partir de las diferencias entre la columna A
        y la columna B, explique qué ocurre con el artículo definido
        \fonol{ʔal}\footnote{El sonido \fon{ʔ} es oclusivo y de articulación
        glotal. El español no cuenta con este sonido en su inventario, pero es
        frecuente en la lengua árabe.} en esta lengua.
  \item Identifique el rasgo que determina la aparición de este fenómeno.
\end{consignas}

\begin{tabular}{@{}ll@{\qquad}ll@{}}
  \toprule
  \multicolumn{2}{@{}l}{\textbf{A}} & \multicolumn{2}{l@{}}{\textbf{B}} \\
  \textit{Transcripción} & \textit{Traducción} & \textit{Transcripción} & \textit{Traducción} \\
  \midrule
  ʔal-qamr           & la luna         & ʔaʃ-ʃams           & el sol \\
  ʔal-walad          & el chico        & ʔat-tiːn           & los higos \\
  ʔal-bint           & la chica        & ʔas-safar          & el viaje \\
  ʔal-ism            & el nombre       & ʔaz-zaffa          & la boda \\
  ʔal-kitaːbul-kabiːr & el gran libro  & ʔan-nahruṭ-ṭawiːl  & el largo río \\
  \bottomrule
\end{tabular}
\end{ejercicio}

\section{La sílaba}\label{sec:03-silaba}

La agrupación de fonemas dentro de una palabra se denomina sílaba. La
\termino{sílaba} es una unidad suprasegmental con estructura jerárquica, dado
que se constituye a partir de un núcleo. El núcleo silábico, que en español es
siempre una vocal, puede presentar elementos que lo preceden y reciben la
denominación de \termino{ataque}. A su vez, puede incluir fonemas posteriores
al núcleo, que se conocen como \termino{coda}. El núcleo y la coda conforman la
\termino{rima}. La estructura completa de la sílaba se visualiza en la
figura~\ref{fig:03-silaba}.

\begin{figure}[htbp]
  \centering
  \begin{forest}
    gramatica, for tree={s sep=4mm}
    [Palabra
      [σ
        [Ataque [Ø]]
        [Rima [Núcleo [a]] [Coda [l]]]]
      [σ
        [Ataque [m]]
        [Rima [Núcleo [i]]]]
      [σ
        [Ataque [d]]
        [Rima [Núcleo [o]] [Coda [n]]]]]
  \end{forest}
  \caption{Estructura interna de la sílaba. Núcleo, ataque y coda de las sílabas
  de \textit{almidón}}\label{fig:03-silaba}
\end{figure}

La palabra \forma{almidón} está formada por tres sílabas: \forma{al -- mi --
dón}. Cada una de esas sílabas muestra una estructura diferente en cuanto a sus
márgenes. En la primera, el núcleo \fonol{a} se acompaña por la coda \fonol{l}.
En la segunda, la sílaba \forma{mi} se forma a partir del ataque \fonol{m}
seguido del núcleo \fonol{i}. La última sílaba está constituida por el núcleo
\fonol{o}, el ataque \fonol{d} y la coda \fonol{n}. Como se puede ver en la
figura~\ref{fig:03-silaba}, las sílabas presentan estructuras diversas. En
español, la sílaba mínima es la que se forma a partir de una única vocal
nuclear, como las primeras sílabas de las palabras \forma{a.mi.go},
\forma{e.le.gir}, \forma{i.le.gal}, \forma{o.la}, \forma{u.so}. El núcleo
vocálico, que constituye la máxima sonoridad de la sílaba, puede estar
acompañado de una semiconsonante que posee menor sonoridad y que se configura,
en términos de Hualde (2013), como un «satélite» de la vocal. Este tipo de
núcleo se observa en el caso de los diptongos de las palabras \forma{ai.re},
\forma{au.la}, \forma{pei.ne}, \forma{pau.ta}.

Asimismo, el ataque o la coda también pueden ser complejos, cuando presentan más
de una consonante, como la primera sílaba de \forma{fra.ter.no},
\forma{trans.por.te} o \forma{cons.tru.ir}. Los grupos consonánticos de ataque
silábico en español son aquellos formados por \fonol{p}, \fonol{b}, \fonol{g},
\fonol{k}, \fonol{f}, \fonol{t} + \fonol{l} o \fonol{r}, por ejemplo
\forma{bro.ma, frá.gil, cro.a.ta}, y el grupo \fonol{dr}, como en
\forma{dra.gón}. Las codas complejas contienen como elemento final siempre
\fonol{s}. Sin embargo, es habitual que las consonantes en coda silábica sufran
procesos de elisión o neutralización.

\begin{ejercicio}\label{ej:03-silabas}
\begin{consignas}
  \item Divida las siguientes palabras en sílabas.

        \forma{barco, implante, halcón, suerte, paseo, protector, antiguo,
        particular, transmutable}
  \item Describa estructuralmente todas las sílabas de las palabras en el
        siguiente cuadro:

        {\footnotesize
        \begin{tabular}{@{}*{9}{c}@{}}
          \toprule
          \multicolumn{3}{c}{Sílaba 3} & \multicolumn{3}{c}{Sílaba 2} & \multicolumn{3}{c}{Sílaba 1} \\
          \cmidrule(r){1-3}\cmidrule(lr){4-6}\cmidrule(l){7-9}
          Ataque & Núcleo & Coda & Ataque & Núcleo & Coda & Ataque & Núcleo & Coda \\
          \midrule
          & & & B & A & R & C & O & \\
          \bottomrule
        \end{tabular}}
  \item Busque palabras que contengan sílabas con las estructuras pedidas abajo
        (V: vocal; C: consonante):
        \begin{exe}
          \ex V: \forma{te.\textbf{a}.tro}
          \ex CV:
          \ex VC:
          \ex CVC:
          \ex CCV:
          \ex CCVC:
          \ex CCVCC:
        \end{exe}
\end{consignas}
\end{ejercicio}

\section{Los rasgos suprasegmentales}\label{sec:03-suprasegmentales}

Además de los segmentos vocálicos y consonánticos que conforman las sílabas, las
palabras y los grupos de palabras presentan \termino{rasgos suprasegmentales}
que se superponen a los segmentos. El acento y la entonación son los rasgos más
importantes para el español. Ambos tienen valor lingüístico tanto a nivel de la
palabra como del sintagma, respectivamente.

\subsection{El acento}

El acento es un rasgo suprasegmental que ocurre dentro del dominio de la sílaba.
Consiste en resaltar una sílaba sobre las otras de su entorno. Esto se realiza
en español a través de distintos recursos acústicos: aumento del tono, la
duración y la energía de pronunciación de la sílaba, entre otras propiedades.
La mayoría de las palabras presentan al menos un acento léxico. En español, la
posición de la sílaba acentuada puede variar.

\begin{itemize}
  \item Las palabras agudas u oxítonas son aquellas que están acentuadas en la
        última sílaba, por ejemplo, \forma{cantar} y \forma{pisó}.
  \item Las palabras graves o paroxítonas son aquellas cuyo acento está en la
        anteúltima sílaba, como \forma{difícil} y \forma{piso}.
  \item Las palabras esdrújulas o proparoxítonas son aquellas con acento en la
        antepenúltima sílaba, como \forma{fantástico} y \forma{cálida}.
\end{itemize}

Si bien cada palabra presenta el acento en una posición, un mismo grupo de
segmentos puede presentar variación en el acento. Este es el caso de las
palabras \fon{ˈteɾmino}, \fon{teɾˈmino} y \fon{teɾmiˈno}. Se puede observar que
la variación de la posición de la sílaba tónica conlleva un cambio de
significado léxico. De esta manera, el acento en español se caracteriza por su
capacidad de distinguir palabras idénticas y, por este motivo, es un rasgo
prosódico que tiene valor fonológico. En resumen, el español es una lengua cuyo
acento es \termino{libre y contrastivo}.

No todas las lenguas presentan esta propiedad. Existen lenguas de \termino{acento
fijo}, en las cuales el acento cae siempre sobre la misma sílaba, como el
francés y el polaco, entre otras.

Si bien de forma aislada todas las palabras presentan un acento primario,
también existen excepciones a la obligatoriedad del acento en español. La mayor
parte de las preposiciones y las palabras funcionales, como las conjunciones o
determinantes, forman un grupo prosódico con la palabra que les sigue; por lo
tanto, no presentan un acento léxico propio.

Asimismo, los adverbios terminados en \forma{-mente} y algunas palabras
compuestas pueden presentar dos acentos. Por ejemplo, la palabra
\forma{fácilmente} posee una sílaba acentuada en la raíz y otra sílaba acentuada
en el sufijo \fon{ˈfasilˈmente}.

\begin{ejercicio}\label{ej:03-acento}
\begin{consignas}
  \item Observe los siguientes corpus del inglés británico y del turco. ¿Qué
        tipo de acento presenta el inglés? ¿Cómo se justifica a partir del
        corpus?
  \item ¿Qué consecuencias trae el cambio en la sílaba tónica en estos pares?
  \item ¿Qué diferencia presenta con el turco? ¿Por qué?
\end{consignas}

\textit{Inglés británico}

\begin{tabular}{@{}lll@{}}
  \fon{ˈʤɜːmən} \sig{alemán}        & \fon{ˈpɾəʊgɾɛs} \sig{el progreso} & \fon{ɪnˈkɾiːs} \sig{aumentar} \\
  \fon{ɪnˈsʌlt} \sig{insultar}      & \fon{ˈtɛlɪfəʊn} \sig{teléfono}    & \fon{ˈɪnkɾiːs} \sig{el aumento} \\
  \fon{ˈmɪnɪstəɾ} \sig{ministro}    & \fon{ˈɪnsʌlt} \sig{el insulto}    & \fon{pɾəʊˈgɾɛs} \sig{progresar} \\
  \fon{ˈfɑːðə} \sig{padre}          & & \\
\end{tabular}

\medskip
\textit{Turco}

\begin{tabular}{@{}lll@{}}
  \fon{baˈkan} \sig{ministro}  & \fon{deˈniz} \sig{mar}          & \fon{teleˈfon} \sig{teléfono} \\
  \fon{aˈçi} \sig{ángulo}      & \fon{baˈba} \sig{padre}         & \fon{meˈraːk} \sig{curiosidad} \\
  \fon{seˈpet} \sig{canasta}   & & \\
\end{tabular}

\begin{resuelto}[Punto (a) resuelto]
El inglés presenta acento libre y contrastivo, como lo muestra el par
\forma{insultar} \fon{ɪnˈsʌlt} y \forma{el insulto} \fon{ˈɪnsʌlt}.
\end{resuelto}
\end{ejercicio}

\subsection{La entonación}

El habla humana transmite información lingüística a través de recursos
acústicos que traspasan el dominio de la palabra. Este fenómeno se conoce como
\termino{entonación} y se produce a partir de variaciones en el tono, la
duración o la velocidad, entre otros.

La entonación posee varias funciones en el lenguaje. Por un lado, permite
distinguir información discursiva al indicar el tipo de enunciado que se
realiza. En este sentido, se vincula con la \termino{modalidad},\footnote{La
modalidad se trata en el \capref{14} de este volumen.} dado que es un recurso
que utiliza el hablante para expresar sus intenciones comunicativas. Por
ejemplo, la secuencia \forma{vas a llegar temprano} se puede interpretar como un
enunciado declarativo \exref{ex:03-modalidad-a}, interrogativo
\exref{ex:03-modalidad-b} o exclamativo \exref{ex:03-modalidad-c}, según el tipo
de entonación que lo acompañe.

\begin{exe}
  \ex\label{ex:03-modalidad}
    \begin{xlist}
      \ex\label{ex:03-modalidad-a} Vas a llegar temprano.
      \ex\label{ex:03-modalidad-b} ¿Vas a llegar temprano?
      \ex\label{ex:03-modalidad-c} ¡Vas a llegar temprano!
    \end{xlist}
\end{exe}

Por otra parte, la entonación permite agrupar palabras en unidades mayores y así
delimitar los enunciados. A través de la presencia o ausencia de pausas o de
cambios tonales, se pueden generar distintos significados para la misma
secuencia lineal de palabras. Así, la entonación, gracias a diversos recursos
acústicos, permite definir o desambiguar las frases que constituyen un
enunciado, fenómeno que se conoce como \termino{fraseo}. La oración \forma{Lupe
dice que cocines} se puede interpretar de dos maneras. Gracias a la presencia de
una pausa, representada con una coma en \exref{ex:03-lupe-a}, el nombre propio
es un vocativo, dado que está usado por el sujeto hablante en función
apelativa, para dirigirse a un interlocutor; mientras que en
\exref{ex:03-lupe-b} el nombre propio funciona como el sujeto del verbo
\forma{decir}.

\begin{exe}
  \ex\label{ex:03-lupe}
    \begin{xlist}
      \ex\label{ex:03-lupe-a} Lupe, dice que cocines. (Alguien se dirige a Lupe
          para pedirle que cocine).
      \ex\label{ex:03-lupe-b} Lupe dice que cocines. (Lupe se dirige a otra
          persona para pedirle que cocine).
    \end{xlist}
\end{exe}

La entonación, además, es el vehículo para codificar información de tipo
pragmático (cortesía, mandato, etc.), sociolectal (dialecto, edad, etc.) o
afectivo (estados de ánimo, humor, etc.).

\begin{recuadro}[El tono]
El tono es un rasgo suprasegmental cuyo correlato acústico es la modulación de
la frecuencia fundamental en la que vibran las cuerdas vocales (F0). La F0 se
calcula sobre la base de la frecuencia de abertura y cierre de las cuerdas
vocales. El tono puede ser alto, medio o bajo, según el rango de la F0. En
español, las variaciones del tono poseen una función pragmática o afectiva; sin
embargo, en otras lenguas llamadas \termino{lenguas tonales}, el tono tiene una
función distintiva a nivel léxico. Por ejemplo, en chino, la sílaba \fonol{ma}
posee diversos significados según su tono: «Con tono alto sostenido significa
\sig{madre}; con tono ascendente empezando desde una altura media,
\sig{cáñamo}; con un contorno bajo ascendente-descendente, \sig{caballo}, y con
tono descendente desde un punto de partida alto, \sig{reñir}» (Hualde, 2013,
p.~260).
\end{recuadro}

\begin{ejercicio}\label{ej:03-puntuacion}
\begin{consignas}
  \item Realice la transcripción ortográfica del siguiente fragmento.
  \item Identifique distintos tipos de enunciados y coloque los signos de
        puntuación correspondientes.
\end{consignas}

\fon{ˈo la ˈko mo es ˈtas te ʃa ˈma βa ˈpa ra kon sul ˈtar te pu ˈðis te kon se
ɣi ˈr el ˈli βro ðe ˈleŋ gwa por fa ˈβor a βi ˈsa me no te ol ˈβi ðes}
\end{ejercicio}

\begin{ejercicio}\label{ej:03-ambiguedad}
Encuentre la ambigüedad que se genera en los siguientes pares de enunciados,
según el tipo de entonación con que se pronuncien. Explique cada caso y
relacione la entonación con el uso de la puntuación.

\begin{exe}
  \ex Cuando llegó, Emilia me entregó el paquete. / Cuando llegó Emilia, me
      entregó el paquete.
  \ex No los aceptaron rápidamente. / No, los aceptaron rápidamente.
  \ex Los clientes que tuvieron problemas deben reclamar en la oficina. / Los
      clientes, que tuvieron problemas, deben reclamar en la oficina.
  \ex Fernando te molestaba por un turno. / Fernando, te molestaba por un
      turno.
\end{exe}

\begin{resuelto}
1. En el primer caso, Emilia es el sujeto del verbo principal
\forma{entregó}, mientras que en el segundo caso es el sujeto del verbo de la
subordinada \forma{llegó}. Si bien es posible interpretar que la persona que
llega y la que entrega el paquete pueden diferir, en el primer caso es más
lógico pensar que las dos acciones son realizadas por la misma persona. En el
primer caso, la frase entonativa termina en \forma{llegó}, mientras que en el
segundo termina en \forma{Emilia}.
\end{resuelto}
\end{ejercicio}

\begin{ejercicio}\label{ej:03-humor}
En estas imágenes se produce un efecto humorístico debido a la ambigüedad en la
interpretación de los rasgos suprasegmentales. Describa qué interpretaciones se
desprenden y cómo se producen los efectos de ambigüedad.

\imagenpendiente{Imágenes A y B del ejercicio}
\end{ejercicio}

\begin{recuadro}[Reflexión final]
En este capítulo se abordaron los contenidos básicos para la comprensión del
aspecto sonoro del lenguaje. En primer lugar, se presentaron las reglas y
excepciones que rigen el sistema de escritura del español; en segundo lugar, se
describieron brevemente los métodos y enfoques para el estudio de los sonidos
del habla, es decir, de la fonética. Finalmente, se expusieron los principios
que explican cómo la organización y modulación del sonido fundamentan la
construcción del significado lingüístico, objeto de estudio de la fonología.
Las herramientas conceptuales y analíticas desarrolladas en este capítulo
constituyen una base indispensable para diversas áreas de aplicación, tales como
la lingüística computacional, los estudios forenses, la lingüística clínica y
la fonoaudiología, entre otras. Con frecuencia no se considera suficientemente
la relevancia de este campo disciplinar. Por ello, es importante destacar que la
fonética y la fonología son saberes transversales que se vinculan no solo con
los demás niveles de la gramática, sino también con múltiples dimensiones del
estudio de la actividad humana.
\end{recuadro}

% --- Secciones finales -------------------------------------------------
\lecturasrecomendadas

\solucionario

\begin{solucion}{\ref{ej:03-viento}}
a. 2 y 3. \forma{por} y \forma{va}: todos los grafemas de la palabra representan
un sonido. 4. \forma{Manu Chao}: hay 7 segmentos; la secuencia \forma{c+h}, en
palabras como \forma{Chao} o \forma{choclo}, representa un solo sonido. 5.
\forma{carretera} \fon{kareˈteɾa}: hay 8 segmentos; la consonante vibrante
múltiple \fonol{r} se escribe en español mediante la grafía doble \forma{rr}
entre vocales. 6. \forma{hombre} \fon{ˈombɾe}: hay 5 segmentos; la letra
\forma{h} en español no representa ningún fonema.

b. 2. \forma{ruta -- carretera -- volverá}. El sistema de representación
ortográfica de las vibrantes puede distinguirse de acuerdo con el contexto de
aparición: la vibrante múltiple \fon{r} se grafica con \forma{rr} entre vocales,
como en \forma{perro}, y con \forma{r} al inicio de palabra o después de las
consonantes \fonol{n}, \fonol{l} o \fonol{s}, como en \forma{rojo},
\forma{Enrique}, \forma{alrededor} y \forma{desratizar}. Por otro lado, la letra
\forma{r} entre vocales corresponde a la vibrante simple \fon{ɾ}, como en
\forma{pero}. 3. \forma{sin -- razón}. En el español de América, los grafemas
\forma{s}, \forma{c} y \forma{z} se pronuncian de la misma manera (seseo). 4.
\forma{hombre -- hambre}. La letra inicial de ambas palabras no representa
ningún sonido.

c. Las palabras \forma{cielo} y \forma{celos} comparten el grafema inicial con
\forma{carretera} y \forma{cuando}. La pronunciación del grafema \forma{c}
delante de las vocales \fon{i}, \fon{e} es similar a la de \forma{sin} y
\forma{razón} en el español americano y canario (comunidades seseantes); en
cambio, delante de \fon{a}, \fon{o} y \fon{u}, \forma{c} se pronuncia como
\fon{k}, como en \forma{casa}, \forma{comida} y \forma{cuenco}.

d. 2. Verdadero; por ejemplo, el grafema \forma{g} se pronuncia \fon{g} en
\forma{gato} y \fon{x} en \forma{gema}.\pendiente{El original dice [ç]} 3.
Falso; por ejemplo, en algunos casos, el grafema \forma{x} representa dos
sonidos, la unión de \fon{k} + \fon{s}, como en \forma{taxi}. 4. Verdadero; por
ejemplo, \forma{rr} en \forma{perro} o \forma{ch} en \forma{chancho}. 5. Falso;
por ejemplo, la letra \forma{h} no representa ningún sonido. 6. Verdadero; por
ejemplo, el grafema \forma{r} puede representar un sonido vibrante múltiple al
inicio de palabra, como en \forma{rata}, y un sonido vibrante simple entre
vocales, como en \forma{para}.

e. 2. \forma{\textbf{g}uante -- \textbf{g}ala -- \textbf{g}erminar --
\textbf{g}ol -- \textbf{g}itano}. Un único grafema puede representar varios
sonidos, como el grafema \forma{g}, que representa el sonido \fon{g} en posición
inicial y delante de \fon{a}, \fon{o} y \fon{u}, pero delante de \fon{e},
\fon{i} se pronuncia como \fon{x}.\pendiente{El original dice [ç]} 3.
\forma{ta\textbf{x}i -- e\textbf{x}amen -- lé\textbf{x}ico -- ne\textbf{x}o}.
Algunos grafemas pueden representar más de un sonido; por ejemplo, el grafema
\forma{x} representa dos sonidos, la unión de \fon{k} + \fon{s}. 4.
\forma{\textbf{ll}amar -- \textbf{y}a -- \textbf{y}erba -- \textbf{ll}eno}. Dos
o más grafemas pueden representar un único sonido; así, en regiones yeístas (ver
recuadro «Seseo y yeísmo»), estos dos grafemas corresponden a la misma
pronunciación, aunque la pronunciación de estas consonantes puede variar en
otras regiones. 5. \forma{\textbf{g}anar -- \textbf{gu}eto -- \textbf{g}ol --
\textbf{gu}iño -- \textbf{g}ula}. Dos grafemas contiguos pueden corresponder a
un único sonido; por ejemplo, la dupla \forma{gu} representa la pronunciación
\fon{g} delante de las vocales anteriores \fon{e}/\fon{i}. 6.
\forma{\textbf{c}osa -- \textbf{c}álido -- \textbf{qu}echua --
\textbf{qu}ímica -- \textbf{c}ubo}. Distintos grafemas representan un solo
sonido; por ejemplo, la dupla \forma{qu} representa la pronunciación \fon{k}
delante de las vocales anteriores \fon{e}/\fon{i}.
\end{solucion}

\begin{solucion}{\ref{ej:03-ortografia}}
a. 1. La ortografía correcta es \forma{tarjetas}. La confusión se produce porque
los grafemas \forma{j} y \forma{g} delante de las vocales \fon{e}/\fon{i} tienen
una pronunciación idéntica. 2 y 4. La ortografía correcta es \forma{venta} y
\forma{favor}. La confusión entre \forma{b} y \forma{v} es habitual ya que ambos
grafemas representan el mismo fonema en todas las variedades del español. 3. La
ortografía correcta es \forma{hielo}. La omisión del grafema \forma{h} se
produce por su falta de manifestación acústica. 5. La ortografía correcta es
\forma{silencio}. 6. La ortografía correcta es \forma{descansen}. Para ciertos
hablantes es habitual confundir los grafemas \forma{s} y \forma{c} delante de
\fonol{e}, \fonol{i}, dado que se pronuncian de igual manera en contextos
seseantes. 7. La ortografía correcta es \forma{tuyo}. 8. La ortografía correcta
es \forma{pollo}. Para los hablantes yeístas es habitual confundir los grafemas
\forma{ll} y \forma{y}, dado que se pronuncian de igual manera. En el caso de
\forma{tuyo} (pronombre posesivo de 2.ª persona singular) y \forma{tullo}
(1.ª persona del verbo \forma{tullir}) se produce un caso de homofonía de dos
palabras que en la escritura diferencian significado.

b. Al igual que ocurre con \forma{tuyo/tullo} y \forma{haya/halla}, otros pares
como \forma{hecho} (verbo \forma{hacer}) / \forma{echo} (verbo \forma{echar}),
\forma{valla} (cerca o vallado) / \forma{vaya} (verbo \forma{ir}), \forma{ay}
(interjección) / \forma{hay} (verbo \forma{haber}) presentan el mismo fenómeno de
homofonía (al menos en la variante rioplatense). La homofonía es un fenómeno en
el cual dos palabras se pronuncian de manera idéntica, pero significan distintas
cosas. En este caso, dado que las palabras se escriben de diferente manera a
pesar de que suenan igual, aunque son homófonas, no son homógrafas.
\end{solucion}

\begin{solucion}{\ref{ej:03-aparato}}
b. Si el velo del paladar está bajo, el sonido que se produce es nasal
(respuesta 2). c. La glotis es el espacio que se forma entre las cuerdas vocales
(respuesta 3). d. Son órganos móviles: lengua, labios, paladar blando, velo del
paladar y úvula. Los alvéolos, el paladar duro y los dientes son articuladores
fijos.
\end{solucion}

\begin{solucion}{\ref{ej:03-vocales}}
\forma{el} \fon{e}: medio cerrado anterior no redondeado. \forma{razón}
\fon{o}: medio cerrado posterior redondeado. \forma{ruta} \fon{u}: cerrado
posterior redondeado. \forma{sin} \fon{i}: cerrado anterior no redondeado.
\end{solucion}

\begin{solucion}{\ref{ej:03-transcripcion}}
B. por -- 7. \fon{poɾ}; C. suerte -- 4. \fon{ˈsweɾte}; D. va -- 3. \fon{ba};
E. Manu Chao -- 5. \fon{manuˈtʃao}; F. viento -- 2. \fon{ˈbjento}; G. el -- 8.
\fon{el}; H. hombre -- 1. \fon{ˈombɾe}; I. carretera -- 9. \fon{kareˈteɾa}.
\end{solucion}

\begin{solucion}{\ref{ej:03-articulacion}}
a.

\begin{tabular}{@{}>{\raggedright\arraybackslash}p{5cm} >{\raggedright\arraybackslash}p{4cm}@{}}
  \toprule
  \textbf{Oclusivos} & \textbf{Fricativos} \\
  \midrule
  \forma{viene; viento; va; carretera; cuando; Babylon; volverá; por} &
  \forma{se; frontera; sin; suerte} \\
  \bottomrule
\end{tabular}

b.

\begin{tabular}{@{}*{5}{Y{1.9cm}}@{}}
  \toprule
  \textbf{bilabial} & \textbf{labiodental} & \textbf{alveolar} & \textbf{palatal} & \textbf{velar} \\
  \midrule
  \fon{ˈmanu} \fon{ˈbjento} \fon{poɾ} & \fon{fɾonˈteɾa} & \fon{raˈson} \fon{ˈsweɾte} & \fon{tʃao} & \fon{ˈkwando} \\
  \bottomrule
\end{tabular}

c. Las palabras \forma{viento} \fon{ˈbjento}, \forma{por} \fon{poɾ} y
\forma{Manu} \fon{ˈmanu} están introducidas por consonantes oclusivas
bilabiales. Las consonantes \fon{p} y \fon{b} contrastan entre sí por la acción
de las cuerdas vocales: \fon{p} es sorda y \fon{b} es sonora. El sonido inicial
de \fon{ˈmanu} es bilabial oclusivo y sonoro. Se diferencia de \fon{b} por la
acción del velo del paladar: \fon{m} es nasal y \fon{b} es oral.

d. Bilabial oclusiva oral sonora: \fon{b} en \fon{ˈbjento}. Alveolar vibrante
múltiple oral sonora: \fon{r} en \fon{ˈruta}. Labiodental fricativa oral sorda:
\fon{f} en \fon{fɾonˈteɾa}. Alveolar lateral oral sonora: \fon{l} en \fon{la}.
Alveolar oclusiva oral sorda: \fon{t} en \fon{kareˈteɾa}.

e. 2. \fon{ɛ}: semiabierta anterior no redondeada.\pendiente{El original dice
«redondeada»} 3. \fon{h}: glotal fricativa oral sorda. 4. \fon{x}: velar
fricativa oral sorda. 5. \fon{ɡ}: velar oclusiva oral sonora. 6. \fon{ʒ}:
postalveolar fricativa oral sonora. 7. \fon{ɔ}: semiabierta posterior
redondeada.

f. 2. \fon{ˈsizmo}: \forma{sismo}. 3. \fon{akoˈplaβa}: \forma{acoplaba}. 4.
\fon{ˈoŋɡo}: \forma{hongo}. 5. \fon{ˈleʧe}: \forma{leche}. 6. \fon{ˈsjero}:
\forma{cierro}. 7. \fon{teˈxaðo}: \forma{tejado}.\pendiente{No figura en la
consigna} 8. \fon{arkeˈoloɣa}: \forma{arqueóloga}. 9. \fon{ˈweso}:
\forma{hueso}. 10. \fon{ˈpeɾa}: \forma{pera}. 11. \fon{ˈkaɲa}: \forma{caña}. 12.
\fon{ˈtaksi}: \forma{taxi}.

g. 2. \forma{caño} \fon{ˈkaɲo}. 3. \forma{llave} \fon{ˈʃaβe}. 4. \forma{dura}
\fon{ˈduɾa}. 5. \forma{neurona} \fon{newˈɾona}. 6. \forma{pienso} \fon{ˈpjenso}.
7. \forma{palo} \fon{ˈpalo}. 8. \forma{hora} \fon{ˈoɾa}. 9. \forma{cereza}
\fon{seˈɾesa}. 10. \forma{puerro} \fon{ˈpwero}. 11. \forma{cama} \fon{ˈkama}.
12. \forma{ley} \fon{lei̯}. 13. \forma{terapeuta} \fon{teɾaˈpewta}. 14.
\forma{jota} \fon{ˈxota}.
\end{solucion}

\begin{solucion}{\ref{ej:03-paresminimos}}
a. Par mínimo presente en el corpus: \fon{ˈombɾe} y \fon{ˈambɾe}. b. Otros
ejemplos de pares mínimos: \fon{mas} \fon{mes}; \fon{ba} \fon{bi}; \fon{ˈmanu}
\fon{ˈmano}; \fon{ˈbjento} \fon{ˈsjento}; \fon{sin} \fon{fin}; \fon{ˈsweɾte}
\fon{ˈmweɾte}; \fon{poɾ} \fon{pon}; \fon{se} \fon{me}.
\end{solucion}

\begin{solucion}{\ref{ej:03-ue}}
a. Se identifican tres pares mínimos en el corpus: \fon{pus} \forma{pus} y
\fon{pes} \forma{pez}, \fon{suɾ} \forma{sur} y \fon{seɾ} \forma{ser},
\fon{ˈluʧe} \forma{luche} y \fon{ˈleʧe} \forma{leche}. b. Dado que la
sustitución de los fonos \fon{e} y \fon{u} en el mismo contexto produce cambio
de significado, como se evidencia en los pares mínimos encontrados, se puede
afirmar que estos fonos son fonemas distintos en español y mantienen una
distribución contrastiva.
\end{solucion}

\begin{solucion}{\ref{ej:03-aleman}}
a. Se puede observar el par mínimo \fon{ˈʃpuːkən} \sig{embrujado} y
\fon{ˈʃpʊkən} \sig{escupir} en el corpus del alemán. b. A partir de la
identificación de este par mínimo se puede concluir que los fonos \fon{uː} y
\fon{ʊ} son fonemas distintos en el sistema fonológico del alemán. Se observa
que su intersustitución genera un cambio de significado; por lo tanto, tienen
una distribución contrastiva en esta lengua.
\end{solucion}

\begin{solucion}{\ref{ej:03-bd}}
a. Si se observa con detenimiento la pronunciación de las consonantes marcadas
en negrita en posición inicial o en medio de la palabra, se puede apreciar que
el contacto de los labios no es idéntico en ambas articulaciones. b. En la
transcripción del corpus se aprecia que después de silencio (es decir, al inicio
de la palabra) los fonemas \fonol{b} y \fonol{d} se producen con cierre completo,
o sea, como oclusivos \fon{b} \fon{d}, mientras que entre vocales ese cierre no
llega a producirse completo y el aire se escapa por entre los articuladores,
produciendo una articulación aproximante \fon{β} \fon{ð}. c. Descripción
articulatoria: \fon{d} alveolar oclusiva oral sonora y \fon{ð} dental
aproximante (fricativa) oral sonora. d. Regla: después de silencio, ocurre el
alófono oclusivo, mientras que en contexto intervocálico la consonante se
realiza como aproximante.
\end{solucion}

\begin{solucion}{\ref{ej:03-sz-espanol}}
b. El corpus no presenta ningún par mínimo. c. En español, \fon{s} y \fon{z}
son alófonos de un mismo fonema y tienen una distribución complementaria, ya que
no comparten contexto. El alófono \fon{s} se realiza después de silencio y entre
vocales, como en \fon{raˈson} y \fon{sin}, y el alófono \fon{z} ocurre delante
de consonante sonora, como en \fon{dezˈðen} y \fon{konsuˈmizmo}.
\end{solucion}

\begin{solucion}{\ref{ej:03-sz-frances}}
b. En francés, los fonos \fon{s} y \fon{z} son fonemas distintos. Ambos ocurren
en los mismos contextos y permiten distinguir significados, es decir que tienen
distribución contrastiva, como se evidencia en los pares mínimos encontrados.
\end{solucion}

\begin{solucion}{\ref{ej:03-quechua}}
a. La vocal \fon{e} es anterior, medio cerrada y no redondeada. La vocal \fon{o}
es posterior, medio cerrada y redondeada. La vocal \fon{u} es posterior, cerrada
y redondeada. b. Los fonos \fon{i} y \fon{e} son alófonos del mismo fonema en
quechua. No comparten contexto y no contrastan significado. El contexto de
aparición de \fon{e} es en contacto con el fonema \fonol{q}, mientras que el
alófono \fon{i} aparece en el resto de los contextos. c. En español, los fonos
\fon{i} y \fon{e} son fonemas distintos, como lo evidencia el par mínimo
\fon{ˈmisa} y \fon{ˈmesa}. Tienen distribución contrastiva y su
intersustitución produce cambio de significado. d. Con los fonos \fon{o} y
\fon{u} del quechua sucede lo mismo que se describió en (b) respecto de \fon{i}
y \fon{e}. e. En español, con los fonos \fon{o} y \fon{u} sucede lo mismo que se
describió en el punto (c) respecto de \fon{i} y \fon{e}, como lo evidencia el
par mínimo \fon{ˈposo} y \fon{ˈpuso}.
\end{solucion}

\begin{solucion}{\ref{ej:03-ruso}}
a. El corpus no presenta ningún par mínimo. b. Ambas vocales son cerradas, pero
se diferencian en la posición que adopta la lengua al producirlas: \fon{i} es
anterior, mientras que \fon{ɨ} es central. c. Estas vocales tienen distribución
complementaria. d. La vocal \fon{i} aparece después de consonantes
palatalizadas (ʲ) y la vocal \fon{ɨ} se produce después de consonantes
velarizadas (ˠ). Palatalización y velarización son procesos fonológicos de
asimilación (ver §\ref{sec:03-procesos}) en los cuales se produce un cambio en
el punto de articulación.\pendiente{El original remite a §8.1.2} Una consonante
palatalizada adelanta hacia el paladar su punto de articulación, mientras que
una consonante velarizada lo desplaza hacia el velo del paladar.
\end{solucion}

\begin{solucion}{\ref{ej:03-bsas}}
b. No hay pares mínimos con contraste de significado. Los dos fonos aparecen en
los mismos contextos, tanto después de silencio, \fon{ˈʃaβe} / \fon{ˈʒaβe}, como
en contexto intervocálico, \fon{ˈpoʃo} / \fon{ˈpoʒo}. c. Dado que no presentan
contraste ni oposición significativa, estos dos fonos son alófonos del mismo
fonema en variación libre; lo que motiva la aparición de uno u otro es el estilo
de habla o las características sociolectales o dialectales de los hablantes que
los producen. d. El hablante 1 corresponde al audio B, ya que produce el sonido
sordo; mientras que el audio A es el hablante 2, que pronuncia la consonante
sonora.
\end{solucion}

\begin{solucion}{\ref{ej:03-catamarca}}
a. No hay pares mínimos con contraste de significado. Las dos consonantes
aparecen en los mismos contextos, tanto después de silencio, \fon{ˈřio} /
\fon{ˈrio}, como en contexto intervocálico, \fon{ˈkařo} / \fon{ˈkaro}. b. La
respuesta es la misma que la del punto (c) del ejercicio~\ref{ej:03-bsas}. c. El
hablante 1 corresponde al audio A, ya que produce una consonante asibilada;
mientras que el audio B es el hablante 2, que pronuncia la vibrante múltiple.
\end{solucion}

\begin{solucion}{\ref{ej:03-video}}
a y b. Es posible que al inicio del video la pronunciación sea más informal o
relajada, mientras que al final se produzca una pronunciación más controlada o
formal. La distribución de ambos alófonos puede estar motivada por el estilo o
el registro. c. Transcripción ortográfica: «Esto es un río. Este es el canal de
la Mariano Moreno, tiene cascada ahora. Miren lo que es la Mariano Moreno, las
calles que vienen de arriba de la Mariano Moreno. La Prado es imposible
transitarla. Miren lo que es\ldots{} esta es la\ldots{} Rojas, miren lo que es
la Rojas». En el video se puede escuchar la alternancia de los alófonos que se
produce durante la grabación, posiblemente debido a la variación del registro.
\end{solucion}

\begin{solucion}{\ref{ej:03-yendo}}
Ciertos errores de ortografía suelen revelar características dialectológicas de
los autores. En este caso, se trata de un hablante que pronuncia la palatal de
inicio de palabra \forma{yendo} como la semiconsonante \fon{j}, lo cual es un
rasgo característico de la pronunciación de la región noroeste de Argentina.
\end{solucion}

\begin{solucion}{\ref{ej:03-griego}}
b. Estos dos sonidos son dos fonemas diferentes y presentan una distribución
contrastiva en español, como evidencia el par mínimo \fon{ˈpaɾa} y \fon{ˈpara}.
c. Sin embargo, hay pares que no muestran contraste, como el caso de
\fon{daˈðoɾ} / \fon{daˈðor} y \fon{akuˈsarte} / \fon{akuˈsaɾte}. Se puede
concluir que en posición final o antes de consonante, las vibrantes se
neutralizan en español. d. En griego, los sonidos \fon{ɾ} y \fon{r} son
alófonos del mismo fonema \fonol{r} y están en distribución complementaria.
Según la descripción fonológica clásica (Kouti, 2010)\pendiente{Kouti 2010 no
está en la bibliografía}, en griego moderno la vibrante se realiza como múltiple
en coda silábica y como simple en otros contextos. e. En español, las
consonantes \fon{ɾ} y \fon{r} son fonemas diferentes, pero en ciertos contextos
no se observa contraste fonémico (hay neutralización). En griego moderno, los
sonidos son alófonos del mismo fonema con distribución complementaria, es
decir, no comparten contexto. f. La representación fonológica de
\forma{acusarte} y de \forma{dador} incluye un archifonema vibrante:
\fonol{akuˈsaRte}, \fonol{daˈdoR}.
\end{solucion}

\begin{solucion}{\ref{ej:03-arabe}}
a. En la columna A, el artículo definido \fonol{ʔal} del árabe estándar
conserva sus rasgos con todas las consonantes iniciales de la palabra que
antecede, mientras que en la columna B la última consonante del artículo
\fonol{l} sufre un proceso de asimilación regresiva. Esto significa que se
asimila a la consonante que le sigue. En este caso, se trata de una asimilación
total, ya que toma todos sus rasgos. b. La asimilación se produce frente a
consonantes coronales, que son aquellas que tienen el punto de articulación en
la zona dentoalveolar, lo que engloba a las consonantes dentales, alveolares y
postalveolares. Con el resto de las consonantes no ocurre este fenómeno, como se
observa en la columna A.
\end{solucion}

\begin{solucion}{\ref{ej:03-silabas}}
a. bar.co -- im.plan.te -- hal.cón -- suer.te -- pa.se.o -- pro.tec.tor --
an.ti.guo -- par.ti.cu.lar -- trans.mu.ta.ble

b.\pendiente{Revisar el orden de las sílabas y algunas celdas del cuadro original}

{\footnotesize
\begin{tabular}{@{}*{9}{c}@{}}
  \toprule
  \multicolumn{3}{c}{Sílaba 3} & \multicolumn{3}{c}{Sílaba 2} & \multicolumn{3}{c}{Sílaba 1} \\
  \cmidrule(r){1-3}\cmidrule(lr){4-6}\cmidrule(l){7-9}
  Ataque & Núcleo & Coda & Ataque & Núcleo & Coda & Ataque & Núcleo & Coda \\
  \midrule
     &    &    & B  & A  & R & C & O  &   \\
     & I  & M  & PL & A  & N & T & E  &   \\
     &    &    &    & HA & L & C & O  & N \\
     &    &    & S  & UE & R & T & E  &   \\
  P  & A  &    & S  & E  &   &   & O  &   \\
  PR & O  &    & T  & E  & C & T & O  & R \\
     & A  & N  & T  & I  &   & G & UO &   \\
  \bottomrule
\end{tabular}

\medskip
\begin{tabular}{@{}*{12}{c}@{}}
  \toprule
  \multicolumn{3}{c}{Sílaba 4} & \multicolumn{3}{c}{Sílaba 3} & \multicolumn{3}{c}{Sílaba 2} & \multicolumn{3}{c}{Sílaba 1} \\
  \cmidrule(r){1-3}\cmidrule(lr){4-6}\cmidrule(lr){7-9}\cmidrule(l){10-12}
  At. & Nú. & Co. & At. & Nú. & Co. & At. & Nú. & Co. & At. & Nú. & Co. \\
  \midrule
  P  & A & R  & T & I &  & C & U &  & L  & A & R \\
  TR & A & NS & M & U &  & T & A &  & BL & E &   \\
  \bottomrule
\end{tabular}}

c. 1. V: \forma{te.\textbf{a}.tro}; 2. CV: \forma{va.\textbf{so}}; 3. VC:
\forma{\textbf{al}.to}; 4. CVC: \forma{\textbf{car}.tu.cho}; 5. CCV:
\forma{\textbf{bro}.ma}; 6. CCVC: \forma{\textbf{plan}.te.o}; 7. CCVCC:
\forma{\textbf{trans}.mi.tir}.
\end{solucion}

\begin{solucion}{\ref{ej:03-acento}}
b. El cambio de acento indica que se trata de dos clases de palabras diferentes:
un sustantivo y un verbo. c. El turco presenta acento fijo. Todas las palabras
del corpus son agudas u oxítonas.
\end{solucion}

\begin{solucion}{\ref{ej:03-puntuacion}}
a y b. Hola, ¿cómo estás? Te llamaba para consultarte, ¿pudiste conseguir el
libro de Lengua? Por favor, avisame. ¡No te olvides!
\end{solucion}

\begin{solucion}{\ref{ej:03-ambiguedad}}
2. \forma{No los aceptaron rápidamente} / \forma{No, los aceptaron
rápidamente}. En el primer caso, \forma{no} niega el verbo \forma{aceptar}, al
formar una única frase con el verbo. En el segundo caso, niega algo dicho
anteriormente y, al estar separado por una pausa, el verbo \forma{aceptar}
mantiene polaridad positiva. 3. En el primer caso, \forma{que tuvieron
problemas} forma una frase entonativa con \forma{clientes} y se interpreta como
un modificador especificativo que restringe el alcance del núcleo y hace
referencia a un grupo específico, «los que tuvieron problemas». En el segundo
caso, al formar una frase entonativa separada, se deben colocar las comas que
representan las pausas o los cambios de tono propios de la oralidad, y
\forma{que tuvieron problemas} se interpreta como un modificador explicativo que
se aplica a todos los clientes (no a un grupo específico) y puede omitirse. 4.
En el primer caso, \forma{Fernando} es el sujeto de la oración. En el segundo,
al separarse con una pausa representada por la coma en la escritura, se
interpreta como un vocativo y constituye una frase entonativa independiente.
\end{solucion}

\begin{solucion}{\ref{ej:03-humor}}
En la foto A, la omisión de la tilde genera un efecto ambiguo. Sin acento
ortográfico no se puede diferenciar en la escritura las formas del verbo en
1.ª persona del presente de indicativo (\forma{mato/asesino}) de las de
3.ª persona del pretérito perfecto simple (\forma{mató/asesinó}). Si bien la
noticia del medio de comunicación intentaba comunicar un hecho pasado («alguien
mató a seis personas ¡y también asesinó a su suegra!»), la omisión de la tilde
permite la lectura que aparece arriba de la foto. La imagen B muestra un chiste
empleado para destacar la importancia de las comas mediante el uso del humor
negro. El cartel contrasta una frase de uso común, \forma{vamos a comer,
niños}, en la cual la palabra \forma{niños} es un vocativo. Al eliminar la coma,
se realiza una lectura sin pausa en la cual \forma{niños} se interpreta como el
COD del verbo \forma{comer}. Esta reestructuración sintáctica, que nace de una
ambigüedad entonativa, resulta humorística.
\end{solucion}

